%%%%%%%%%%%%%%%%%%%%%%%%%%%%%%%%%%%%%%%%%%%%%%%%%%%%%%%%%
% Niniejszy plik przedstawia przykładowy skład 
% pracy dyplomowej na Wydziale Matematyki PWr. 
% 
% Autorzy: 
% Damian Fafuła
% Michał Kijaczko
% Jakub Michalczak
% Maciej Miśta
% Dagmara Nowak
% Tomasz Skalski
% Wojciech Słomian
%
%% Data utworzenia: 8.05.2018
% Numer wersji: 1
%
% Poniższą formatkę można rozpowszechniać i edytować 
% pod warunkiem zachowania numeru wersji, 
% informacji o autorach i dodaniu informacji 
% o wprowadzonych zmianach.
%
%%%%%%%%%%%%%%%%%%%%%%%%%%%%%%%%%%%%%%%%%%%%%%%%%%%%%%%%%
% Domyślną opcją jest: praca magisterska, język polski.
% W przypadku pracy pisanej w języku angielskim dodajemy 
% opcję [english].
% Dla pracy licencjackiej dodajemy opcję [licencjacka].
% Dla pracy inżynierskiej dodajemy opcję [inzynierska].
% Dopuszczalne są podwójne opcje, np. [licencjacka, english].
% Opcje dodajemy w kwadratowym nawiasie przy \documentclass.
%
%
%%%%%%%%%%%%%%%%%%%%%%%%%%%%%%%%%%%%%%%%%%%%%%%%%%%%%%%%%
\documentclass[licencjacka]{pwr_wmat_praca_dyplomowa}

\usepackage{enumitem}
\usepackage{xcolor}
\usepackage{mathtools}
%%%%%%%%%%%%%%%%%%%%%%%%%%%%%%%%%%%%%%%%%%%%%%%%%%%%%%%%%
%              DANE DO PRACY
%
% W przypadku pracy dyplomowej w języku angielskim nie jest konieczne 
% wypełnianie pól: \tytul{}, \kierunek{}, \specjalnosc{}, 
%                  \streszczenie{}, \slowakluczowe{}.
%%%%%%%%%%%%%%%%%%%%%%%%%%%%%%%%%%%%%%%%%%%%%%%%%%%%%%%%%
%
% Imię i nazwisko autora
\autor{Imię i nazwisko dyplomanta}
%
% Tytuł pracy dyplomowej 
\tytul{Tytuł pracy dyplomowej} 
\tytulang{Tytuł pracy dyplomowej w języku angielskim}
%
% Tytuł / stopień / imię i nazwisko opiekuna
\opiekun{dr inż. Dawid Huczek}
%
% Kierunek studiów wybieramy spośród następujących:
% 1) Matematyka
% 2) Matematyka i Statystyka
% 3) Matematyka stosowana
\kierunekstudiow{Matematyka}
%
% Kierunek studiów po angielsku wybieramy spośród następujących:
% 1) Mathematics
% 2) Mathematics and Statistics
% 3) Applied Mathematics
\kierunekstudiowang{Mathematics}
%
% Specjalność wybieramy spośród następujących: 
% KIERUNEK: Matematyka
% 1) Matematyka teoretyczna,
% 2) Statystyka matematyczna,
% 3) Matematyka finansowa i ubezpieczeniowa,
%
% KIERUNEK: Matematyka i Statystyka
% 4) Matematyka,
% 5) Statystyka i analiza danych, 
%
% 6) -- (w przypadku braku specjalizacji).
\specjalnosc{Matematyka teoretyczna} 
%
% Specjalność w języku angielskim wybieramy spośród następujących:
% KIERUNEK: Matematyka
% 1) Theoretical Mathematics,
% 2) Mathematical Statistics,
% 3) Financial and Actuarial Mathematics,
%
% KIERUNEK: Matematyka i Statystyka
% 4) Mathematics,
% 5) Statistics and Data Analysis,
%
% KIERUNEK: Applied Mathematics
% 6) Financial and Actuarial Mathematics, 
% 7) Mathematics for Industry and Commerce,
% 8) Computational Mathematics,
% 9) Modelling, Simulation and Optimization.
%
% 10) -- (w przypadku braku specjalizacji).
\specjalnoscang{Theoretical Mathematics} 
%
% Krótkie streszczenia po polsku i angielsku
% - nie dłuższe niż 530 znaków.
\streszczenie{Tutaj piszemy krótkie streszczenie pracy (nie powinno być dłuższe niż 530 znaków).}
\streszczenieang{Tutaj piszemy krótkie streszczenie pracy w języku angielskim (nie powinno być dłuższe niż 530 znaków).}
%
% Podajemy najważniejsze słowa kluczowe po polsku i angielsku
% - w obu przypadkach, nie więcej niż 150 znaków.
\slowakluczowe{tutaj podajemy najważniejsze słowa kluczowe (łącznie nie powinny być dłuższe niż 150 znaków).}  
\slowakluczoweang{tutaj podajemy najważniejsze słowa kluczowe w języku angielskim (łącznie nie powinny być dłuższe niż 150 znaków)}
%
%
%%%%%%%%%%%%%%%%%%%%%%%%%%%%%%%%%%%%%%%%%%%%%%%%%%%%%%%%%
% Definicje, lematy, twierdzenia, przykłady i wnioski
% Komendy wywołujące twierdzenia, definicje, itd., 
% czyli 'theorem', 'definition', 'corollary', itd., 
% można zmienić wedle uznania.
\theoremstyle{plain}
\newtheorem{theorem}{Twierdzenie}
\numberwithin{theorem}{chapter}
\newtheorem{lemma}[theorem]{Lemat} 
\newtheorem{corollary}[theorem]{Wniosek}
\newtheorem{fact}[theorem]{Fakt}
\theoremstyle{definition}
\numberwithin{theorem}{chapter}
\newtheorem{definition}[theorem]{Definicja} 
\newtheorem{example}[theorem]{Przykład}
\newtheorem{note}[theorem]{Uwaga}
%%%%%%%%%%%%%%%%%%%%%%%%%%%%%%%%%%%%%%%%%%%%%%%%%%%%%%%%%


%%%%%%%%%%%%%%%%%%%%%%%%%%%%%%%%%%%%%%%%%%%%%%%%%%%%%%%%%
%%%%%%%%%%%%%%%%%%%%%%%%%%%%%%%%%%%%%%%%%%%%%%%%%%%%%%%%%
\begin{document}
\frontmatter
\maketitle
\mainmatter
\tableofcontents
%\listoffigures
%\listoftables



%\chapter{Rozdział pierwszy}
%Tabela \ref{tab:przykladowa} przedstawia przykładową tabelę. Do tworzenia tabeli służą m.in. środowiska \texttt{tabular} oraz \texttt{table}. Istnieje możliwość numeracji dwustopniowej, gdzie pierwsza cyfra oznacza numer rozdziału, a druga – kolejny numer tabeli w tym rozdziale. Tytuł powinien znajdować się centralnie nad tabelą, $12$ pkt odstępu od tekstu zasadniczego nad i pod tabelą wraz z tytułem. Jeśli tabela jest cytowana – należy podać centralnie pod tabelą źródło jej pochodzenia, np. opracowanie własne, opracowano na podstawie danych z GUS.
%\begin{table}[ht]
%\caption{Podstawowa Tabela}
%\centering
%\begin{tabular}{ccc}
%\hline
%\hline                       
%Państwo & PKB (w milionach USD )& Stopa bezrobocia  \\  [0.5ex] 
%\hline 
%Stany Zjednoczone & 75 278 049 & 4,60\%  \\
%Chiny & 11 218 281 & 4,10\%   \\
%Japonia & 4 938 644 & 3,10\%  \\
%Niemcy & 3 466 639 & 6,00\%   \\
%Wielka Brytania & 2 629 188 & 4,60\%  \\ [1ex]  
%\hline 
%\end{tabular}
%\caption*{\textit{Źródło: opracowanie własne}}
%\label{tab:przykladowa} 
%\end{table}
%
%Do cytowania używamy komendy \texttt{cite}. W nawiasie klamrowym podajemy klucz, którego użyliśmy w pliku \emph{bibliografia.bib}. Przykład: \cite{einstein} lub \cite[chap. 2]{latexcompanion}.
%
%\section{Podrozdział pierwszy}
%
%\begin{table}[H]
%\caption{Podstawowa Tabela}
%\centering
%\begin{tabular}{ccc}
%\hline
%\hline                       
%Państwo & PKB (w milionach USD )& Stopa bezrobocia  \\  [0.5ex] 
%\hline 
%Stany Zjednoczone & 75 278 049 & 4,60\%  \\
%Chiny & 11 218 281 & 4,10\%   \\
%Japonia & 4 938 644 & 3,10\%  \\
%Niemcy & 3 466 639 & 6,00\%   \\
%Wielka Brytania & 2 629 188 & 4,60\%  \\ [1ex]  
%\hline 
%\end{tabular}
%\caption*{\textit{Źródło: opracowanie własne}}
%\label{tab:przykladowa2} 
%\end{table}
%
%\section{Podrozdział drugi}
%
%Rysunki do pracy dyplomowej należy wstawiać w sposób podobny do wstawiania tabel, z~zasadniczą różnicą polegającą na tym, że podpis powinno umieszczać się centralnie pod rysunkiem, a nie powyżej niego. Numeracja i sposób cytowania pozostają bez zmian, przy czym tabele i rysunki nie mają numeracji wspólnej, np. po Tabeli \ref{tab:przykladowa2} występuje Rysunek \ref{rys1} (o ile jest to pierwszy rysunek rozdziału pierwszego), a nie Rysunek $1.3$.
%
%\begin{figure}[ht]
%
%\centering
%                     
%\includegraphics[scale=0.27]{logo_w13.jpg}
%\caption{Podstawowy Rysunek}\label{rys1}
%\end{figure}
%\label{rys:przykladowy} 
%
%
%
%\chapter{Definicje, lematy, twierdzenia, przykłady i wnioski}
%Definicje, lematy, twierdzenia, przykłady i wnioski piszemy w pracy tak:
%\begin{definition}[Martyngał]
%Tu piszemy treść definicji martyngału.
%\end{definition}
%\begin{lemma}[]% w nawiasie kwadratowym można napisać jego nazwę
%Tu piszemy treść lematu.
%\end{lemma}


% MOJE ================================================
{\backmatter \chapter{Wstęp}}
\section{Motywacje}
Studia matematyczne podjąłem w dość późnym wieku, już po ukończeniu innych kierunków. Od zawsze jednak matematyka pociągała mnie jako klucz do zrozumienia pozostałych dziedzin, jako narzędzie pozwalające dostrzec lepiej i głębiej urodę przyrody. Wykorzystując matematykę, możemy uchwycić reguły, wzorce i schematy które opisują niewyobrażalną złożoność świata w którym żyjemy. Matematykę postrzegam jako klamrę spinające wszystkie dziedziny intelektualnej działalności człowieka, począwszy od fizyki, poprzez chemię, biologię, aż po psychologię czy nauki społeczne. Decyzję o rozpoczęciu studiów matematycznych podjąłem w lecie siedząc nad rzeką i przyglądając się hipnotyzującym zawirowaniom na powierzchni wody płynącej w rzece Odrze przepływającej przez Wrocław, oraz leniwemu ruchowi fraktalnych kształtów chmur na błękitnym niebie. Byłem wówczas absolwentem wyższej uczelni, w którym wciąż żyło pragnienie poznawania głębiej wszystkich tych pięknych zjawisk, które nas otaczają. Byłem również świeżo po lekturze inspirującej popularnonaukowej książki Jamesa Gleicka zatytułowanej ,,Chaos'' \cite{gleick1987chaos}...
\section{Historia badań nad układami dynamicznymi}
\subsection{Od starożytności do Newtona}
\subsection{Problem stabilności układu słonecznego}
\subsection{Lorenz i jego atraktor}
\subsection{Mandelbrodt i fraktale}
\subsection{Chaos}
Li-Yorke, Von Neumann, Birkhoff, Smale, Szarkowski, Devaney






\chapter{Wprowadzenie}

\section{Definicje}



\begin{enumerate}
\item ciąg
\item podciąg
TODO: Dopisać do sekcji definicji definicję punktu skupienia ciągu.
\item ciąg Cauchy'ego (pojawia się w definicji zupełności, TODO ma być tutaj czy w faktach ogólnych dotyczących zwartych przestrzeni metrycznych?)


\item topologia zbieżności jednostajnej (pojawia się w definicji metryki na przestrzeni odwzorowań trójkątnych)
%\item metryka jest
%\item przestrzeń metryczna jest
%\item kula otwarta jest
%\item zbiór otwarty (korzystamy z pojecia kuli otwartej) jest
%\item pokazac ze kula otwarta jest zbiorem otwartym jest
%\item otoczenie punktu jest
%\item sąsiedztwo punktu jest
%\item wnętrze zbioru jest 
%\item domknięcie jest
%\item zbiór domknięty w przestrzeni metrycznej, jest
%\item zbiór brzegowy, jest
\item kula domknięta






%\item gęstość w przestrzeni metrycznej jest
\item zbiór nigdziegęsty
%\item pokrycie zbiorami otwartymi jest
%\item przestrzen metryczna zwarta, jest
%\item ośrodkowość
\item \textcolor{red}{baza topologii w przestrzeni zwartej (kule o środkach w ośrodku i promieniach wymiernych - potrzebne w dowodzie ogólnym patrz TODO na czerwono)}
%\item zupelnosc, jest
\item punkt izolowany w przestrzeni metrycznej


\item Ciągłość funkcji na przestrzeni metrycznej 
\item (zbiór $C(X)$)


%\item uklad dynamiczny, jest
%\item orbita, jest
%\item iterata odwzorowania, jest
\item iteracja odwzorowania
%\item punkt stały, jest
%\item punkt okresowy odwzorowania, jest
\item orbita okresowa
%\item niezmienniczosc zbioru ze wzgledu na odwzorowanie, jest


%\item entropia topologiczna Adler, jest
%\item entropia topologiczna Bowen, jest
\item entropia topologiczna - dowod rownowaznosci (ksiazka misiurewicz combinatorial dynamics and entropy in dimension one oraz ksiazka ruette rozdzial 4)

%\item topologiczna tranzytywnosc, jest
%\item wrazliwosc na warunki poczatkowe, jest
%\item RODZAJE CHAOSU (Devaneya, jest
%\item chaos Li Yorka), jest


\item napisac ze chaotycznosc odwzorowania rozumiem przez chaotycznosc odpowiedniego ukladu dynamicznego

\item zbiór rezydualny
\item separable, second category (czyli 1 i 2 kategoria bairea)
\item twierdzenie baire'a dla przestrzeni metrycznych zupełnych napisac
\item g-delta




\item odwzorowania trójkątne, zbior $C_\triangle(X \times I)$

\end{enumerate}

% Przyklad cytowania z numerem strony
%\begin{definition}[Metryka]
%TODO
%\cite[s.~103]{kuratowski1977wstep}
%\end{definition}


\subsection{Definicje i fakty ogólne}

\begin{definition}[Ciąg]
Ciąg $(x_n)_{n=1}^\infty$ jest to funkcja określona na zbiorze liczb naturalnych. Wartości tej funkcji dla kolejnych liczb naturalnych nazywamy wyrazami ciągu i oznaczamy: $x_1, x_2, \ldots$. Niekiedy ciąg oznaczamy skrótowo przez $(x_n)$. W pracy rozważać będziemy ciągi w przestrzeniach metrycznych, czyli funkcje postaci $x: \mathbb{N} \rightarrow X,$ gdzie $X$ jest rozważaną przestrzenią metryczną

\end{definition}

\begin{definition}[Podciąg]
TODO
\end{definition}

\begin{definition}[Granica ciągu]
TODO
\end{definition}

TODO: granica górna granica dolna - występuje przy entropii Bowena

\begin{definition}[Punkt skupienia ciągu]
TODO
\end{definition}

\begin{definition}[Ciąg Cauchy'ego]
Ciągiem Cauchy'ego w przestrzeni metrycznej $(X, \rho)$ nazywamy ciąg $(x_n)$ spełniający warunek:
$$\forall_{\epsilon>0} \, \exists_{N \in \mathbb{N}} \, \forall_{n,m > N} \, \rho(x_n, x_m) < \epsilon.$$
\end{definition}




\subsection{Definicje dotyczące zwartych przestrzeni metrycznych}

\begin{definition}[Metryka]
Metryką na zbiorze $X$ nazywamy funkcję $\rho : X \times X \longrightarrow \mathbb{R}_+ \cup \{0\}$ spełniającą następujące warunki:
\begin{enumerate}
\item $\forall_{x,y \in X}: \rho(x,y)=0 \iff x=y$,
\item $\forall_{x,y \in X}: \rho(x,y) = \rho(y,x)$,
\item \label{nierownosc_trojkata} $\forall_{x,y,z \in X}: \rho(x,y) \leq \rho(x,z) + \rho(z,y)$.
\end{enumerate}
Warunek \ref{nierownosc_trojkata} nazywany jest zwykle nierównością trójkąta.

\end{definition}

\begin{definition}[Przestrzeń metryczna]
Przestrzenią metryczną nazywamy parę $(X, \rho)$, gdzie $X$ jest zbiorem a $\rho$ zdefiniowaną na nim metryką. Czasami, tam gdzie nie będzie prowadziło to do nieporozumień, przestrzeń metryczną $(X, \rho)$ będziemy oznaczać przez samo $X$.
\end{definition}

\begin{definition}[Kula otwarta]
Kulą otwartą w przestrzeni metrycznej $(X, \rho)$ nazywamy zbiór: $K(s, r) = \{x \in X: \rho(s, x) < r\}$.
Punkt $s$ nazywamy wówczas środkiem kuli $K$, a $r \in \mathbb{R}_+ \cup \{0\}$ jej promieniem.
\end{definition}

\begin{definition}[Zbiór otwarty]
Zbiór $A$ w przestrzeni metrycznej $(X, \rho)$ nazywamy otwartym wtedy i tylko wtedy gdy, dla każdego $a \in A$ istnieje kula otwarta o środku w $a$, zawierająca się w $A$.
\end{definition}

\begin{lemma}[Kula otwarta w przestrzeni metrycznej jest zbiorem otwartym]\label{kula_otwarta_jest_zbiorem_otwartym}
W przestrzeni metrycznej $(X, \rho)$, każda kula otwarta jest zbiorem otwartym.
\end{lemma}

\begin{proof}
Niech $(X, \rho)$ będzie przestrzenią metryczną. Weźmy dowolną kulę otwartą $K = K(s, r) \subseteq X$. Weźmy teraz dowolny punkt $k \in K$. Wiemy, że $\rho(s, k) < r$, czyli $r - \rho(s, k) = \epsilon > 0$. Zatem, jeżeli weźmiemy $r_k = \frac{\epsilon}{2}$ to $K_2 = K(k, r_k) \subset K$, ponieważ 

\begin{equation}
\begin{aligned}
\forall_{a \in K_2} \\
& \rho(s, a) \leq \rho(s, k) + \rho(k, a) \leq \rho(s, k) + r_k \\
& = \rho(s, k) + \frac{\epsilon}{2} = \rho(s, k) + \frac{r}{2} - \frac{\rho(s, k)}{2} \\
& = \frac{\rho(s, k)}{2} + \frac{r}{2} < r,
\end{aligned}
\end{equation}
czyli $\forall_{a \in K_2} \, a \in K.$ Z dowolności $k$ otrzymujemy, że kula $K$ jest otwarta a z dowolności wyboru kuli $K$ otrzymujemy tezę.
\end{proof}


\begin{definition}[Otoczenie punktu]
W przestrzeni metrycznej $X$ zbiór $V \subseteq X$ nazywamy otoczeniem punktu $x$, jeżeli istnieje zbiór otwarty $U \subseteq X,$ taki że $x \in U \subseteq V$.
\end{definition}

\begin{definition}[Otoczenie zbioru]
W przestrzeni metrycznej $X$ zbiór $V \subseteq X$ nazywamy otoczeniem zbioru $A$, jeżeli istnieje zbiór otwarty $U \subseteq X,$ taki że $A \subseteq U \subseteq V$.
\end{definition}

%\begin{definition}[Sąsiedztwo punktu]
%Sąsiedztwem punktu $x \in X$ nazywamy każdy zbiór $V_x = V \setminus \{x\}$, gdzie $V$ jest pewnym otoczeniem punktu $x$.
%\end{definition}


\begin{definition}[Wnętrze zbioru]
W przestrzeni metrycznej, wnętrzem zbioru $A$ nazywamy zbiór wszystkich punktów, które należą do $A$ wraz z pewnym swoim otoczeniem.
\end{definition}

\begin{definition}[Domknięcie zbioru]
Domknięciem zbioru $A \subset X$ nazywamy zbiór $\{x \in X: \forall_{r>0} \, K(x, r) \cap A \neq \emptyset\}$. Domknięcie zbioru $A$ oznaczamy przez $\bar{A}$ lub $\textrm{Cl}(A)$.
\end{definition}

\begin{definition}[Zbiór domknięty]
Zbiór $A$ jest domknięty gdy $A = \textrm{Cl}(A)$.
\end{definition}

\begin{definition}[Zbiór gęsty]
Dla danej przestrzeni metrycznej $(X, \rho)$. Zbiór $A \subset X$ nazwiemy gęstym, gdy $\forall_{x \in X} \, \forall_{\epsilon>0} \, \exists_{a \in A} : \rho(x, a) < \epsilon$.
\end{definition}

\begin{definition}[Zbiór brzegowy]
Zbiór $A$ nazwiemy brzegowym, gdy ma puste wnętrze.
\end{definition}

\begin{definition}[Zbiór nigdziegęsty]
Dla danej przestrzeni metrycznej $(X, \rho)$. Zbiór $A \subset X$ nazwiemy nigdziegęstym, gdy wnętrze domknięcia tego zbioru jest puste.
\end{definition}

\begin{definition}[Pokrycie otwarte]
\cite[s.~195]{ruette2017chaos}
Pokryciem otwartym przestrzeni metrycznej $X$ nazywamy rodzinę zbiorów otwartych $(U_i)_{i \in I}$ taką, że $X=\bigcup_{i \in I}U_i$.
\end{definition}

\begin{definition}[Przestrzeń metryczna zwarta, definicja pokryciowa]
Na podstawie \cite[s.~196]{ruette2017chaos}.
Przestrzeń metryczną $(X, \rho)$ nazywamy zwartą jeżeli każde pokrycie otwarte $(U_i)_{i \in I}$ tej przestrzeni zawiera podpokrycie skończone, to znaczy istnieje skończony zbiór indeksów $J \subset I,$ taki że $X = \bigcup_{i \in J}U_i$.
Podzbiór $Y \subset X$ jest zwarty jeżeli $(Y, \rho)$ jest przestrzenią zwartą.
\end{definition}

\begin{theorem}[Ciągowa charakteryzacja zwartości przestrzeni metrycznej]
\label{rownowaznosc_definicji_pokryciowej_i_ciagowej}
Przestrzeń metryczna $(X, \rho)$ jest zwarta wtedy i tylko wtedy, gdy z każdego ciągu $(x_n) \subset X$ można wybrać podciąg zbieżny.
\end{theorem}

\begin{proof}
\textcolor{red}{TODO: przyjzec sie temu dowodowi czy pasuje teraz do sformulowania twierdzenia oraz sprobowac poprawic dowod  implikacji w lewa strone (problem z tym ze rodzina tylko przeliczalna a nie uwzgledniamy nieprzeliczalnych)}
$(\Rightarrow)$
Załóżmy, że z dowolnego pokrycia $X$ zbiorami otwartymi możemy wybrać podpokrycie skończone. Weźmy dowolny ciąg $(x_n) \in X$. Pokażę, że ten ciąg ma punkt skupienia w $X$. Załóżmy, że ciąg $(x_n)$ nie ma punktu skupienia w $X$, to znaczy $\forall_{x \in X}$ $x$ nie jest punktem skupienia $x_n$, czyli $\forall_{x \in X} \, \exists_{\epsilon > 0}  K(x, \epsilon) \setminus \{x\} \cap \{x_n\}_{n=1}^{\infty} = \emptyset$. Utwórzmy pokrycie $X$ w następujący sposób: $$\mathcal{A} = \bigcup_{x \in X} K(x, \epsilon_x),$$ gdzie $\epsilon_x > 0$ jest liczbą spełniającą warunek $K(x, \epsilon_x) \setminus \{x\} \cap \{x_n\}_{n=1}^{\infty} = \emptyset$. Oczywiście rodzina $\mathcal{A}$ jest pokryciem otwartym przestrzeni $X$. Zgodnie z założeniem, z pokrycia $\mathcal{A}$ możemy wybrać podpokrycie skończone. Każdy zbiór $A \in \mathcal{A}$ zawiera najwyżej jeden element ciągu $(x_n)$. Skończenie wiele zbiorów zawierających po co najwyżej jednym elemencie ciągu prowadzi do wniosku, że przestrzeń $X$ zawiera skończenie wiele elementów nieskończonego ciągu co stanowi sprzeczność.

Oznaczmy przez $g$ punkt skupienia ciągu $(x_n)$. Oznaczmy przez $y_1$ pierwszy element ciągu $(x_n)$ różny od $g$. przez $y_2$ pierwszy element ciągu $(x_n)$ różny od $g$ i taki, że $y_2 \in K(g, \frac{d}{2})$, gdzie $d = \rho(y_1, g)$.
Następnie niech $y_n$ będzie pierwszym elementem ciągu $(x_n)$ różnym od $g$ i spełniającym warunek $y_n \in K(g, \frac{d_{n-1}}{2})$, gdzie $d_{n-1} = \rho(y_{n-1}, g)$. Powstały w ten sposób podciąg $(y_n) \subset (x_n)$ jest zbieżny do $g \in X$. W przypadku gdyby na którymś etapie powyższej konstrukcji nie można było wskazać kolejnego elementu różnego od $g$, znaczyłoby to, że od pewnego miejsca ciąg $(x_n)$ jest stale równy $g$, a więc zbieżny.

$(\Leftarrow)$
Złóżmy, że $X$ spełnia definicję pokryciową.
Dowód przez kontrapozycję. Pokażemy, że jeżeli istnieje pokrycie $X$, z którego nie można wybrać podpokrycia skończonego to istnieje ciąg z którego nie można wybrać podciągu zbieżnego.

Załóżmy zatem, że nieskończona, przeliczalna rodzina zbiorów otwartych $\{U_i\}_{i=1}^{\infty}$ stanowi pokrycie X. Skonstruujmy następujący ciąg: 
$$x_1 \notin U_1$$
$$x_n \notin \bigcup_{k=1}^{n-1} U_k$$ Zawsze znajdziemy taki $x_n \in X$, gdyż gdyby taki nie istniał znaczyłoby to, że $\bigcup_{i=1}^{n-1}{U_{k_i}} = X$, czyli istniałoby podpokrycie skończone. Załóżmy teraz, że ciąg $(x_n)$ ma podciąg zbieżny. Wtedy istniałaby jego granica, tj. element $g \in X,$ taki że dla każdego $\epsilon > 0$ nieskończenie wiele elementów ciągu leży wewnątrz kuli $B(g, \epsilon)$, z kolei z otwartości zbiorów $U_n$ możemy wybrać taki $\epsilon$, że $B(g, \epsilon) \subset U_k$, gdzie $U_k$ jest dowolnym zbiorem z pokrycia X, zawierającym g. Z konstrukcji naszego ciągu wynika jednak, że $\forall_{n \in N} \forall_{m>n} \, x_m \notin U_n$, czyli do każdego $U_n$ należy jedynie skończenie wiele elementów. Zatem otrzymujemy sprzeczność. Czyli ciąg $(x_n)$ nie zawiera podciągu zbieżnego.
\end{proof}


\subsection{Fakty ogólne, dotyczące zwartych przestrzeni metrycznych}


\begin{definition}[Ośrodkowość]
Przestrzeń metryczna $(X, \rho)$ jest ośrodkowa, gdy zawiera podzbiór przeliczalny i gęsty.
\end{definition}


\begin{definition}[Zupełność]
Przestrzeń metryczna jest zupełna jeśli każdy ciąg Cauchy'ego elementów tej przestrzeni jest zbieżny do elementu tej przestrzeni.
\end{definition}


\begin{theorem}[Ośrodkowość zwartych przestrzeni metrycznych] \label{zwarta_jest_osrodkowa}
Jeżeli przestrzeń metryczna jest zwarta to jest również ośrodkowa.
\end{theorem}

\begin{proof}
Niech $(X, \rho)$ będzie przestrzenią metryczną zwartą. Oznacza to, że z każdego pokrycia $X$ zbiorami otwartymi, można wybrać podpokrycie skończone.

Dla $i \in \mathbb{N}$, niech $\mathcal{A}_{i} = \{ K(x, \frac{1}{i}) : x \in X \}$, $\mathcal{A}_{i}$ jest oczywiście pokryciem $X$ zbiorami otwartymi, zatem z każdej rodziny $\mathcal{A}_{i}$ możemy wybrać podpokrycie skończone, oznaczmy je przez $\mathcal{B}_{i}$. Rozważmy teraz następującą rodzinę: 
$$B = \bigcup\limits_{i=1}^{\infty} \mathcal{B}_{i}.$$
Składa się ona ze zbiorów postaci $K(x, \frac{1}{i})$ dla pewnych $i \in \mathbb{N}$ i $x \in X$ i jako przeliczalna suma rodzin skończonych jest przeliczalną rodziną zbiorów.
Niech teraz 
$$C = \{ s : \exists_{i \in \mathbb{N}} \, K\left(s, \frac{1}{i}\right) \in B \}.$$
Zbiór $C$ jest równoliczny ze zbiorem $B$, a więc przeliczalny.
Ponadto $C$ jest gęsty w $X$:
Niech $\epsilon > 0$. Weźmy dowolny $x \in X$. Oczywiście istnieje $i_0 \in \mathbb{N},$ takie że $\frac{1}{i_0} < \epsilon$. Wiemy, że $\mathcal{B}_{i_0}$ stanowi pokrycie $X$ i składa się ze zbiorów postaci $K(s, \frac{1}{i_0})$, gdzie $s \in X$. Skoro $\mathcal{B}_{i_0}$ jest pokryciem $X$ to $\exists_{K(s, \frac{1}{i_0}) \in \mathcal{B}_{i_0}} : x \in K(s, \frac{1}{i_0})$, czyli $\rho(x, s) < \frac{1}{i_0} < \epsilon$, natomiast $s \in C$ z definicji zbioru $C$

Czyli X zawiera podzbiór przeliczalny i gęsty, zatem przestrzeń $(X, \rho)$ jest ośrodkowa.
\end{proof}


\begin{theorem}[Zupełność zwartych przestrzeni metrycznych] \label{zwarta_jest_zupelna}
Każda zwarta przestrzeń metryczna jest zupełna.
\end{theorem}

\begin{proof}
Niech $(X, \rho)$  będzie przestrzenią zwartą.
Weźmy ciąg Cauchy'ego $(x_n)$ elementów $X$. Ze zwartości (twierdzenie \ref{rownowaznosc_definicji_pokryciowej_i_ciagowej}) wynika, że istnieje podciąg $(y_{n_k})$ ciągu $(x_n)$ zbieżny do jakiegoś $g \in X$. Zatem $\forall_{\epsilon > 0} \, \exists_{n_0 \in \mathbb{N}} \, \forall_{n > n_0} \, \rho(y_n, g) < \epsilon$.
Pokażemy, że $(x_n)$ zbiega do $g$.

Ustalmy $\epsilon > 0$. Z faktu, że $(x_n)$ jest ciągiem Cauchy'ego wynika, że $$\exists_{N \in \mathbb{N}} \, \forall_{n, m > N} \, \rho(x_m, x_n) < \frac{\epsilon}{2}$$ Ponadto $$\exists_{n_0 \in \mathbb{N}} \, \forall_{n > n_0} \, \rho(y_n, g) < \frac{\epsilon}{2}$$ Weźmy teraz element $y_{k}$ podciągu $(y_n)$, taki że $y_{k} = x_M$, gdzie $M > N \land M > n_0$. Wówczas 
$$M > N \implies \forall_{n > N} \, \rho(x_M, x_n) < \frac{\epsilon}{2}$$
$$M > n_0 \implies \rho(x_M, g) < \frac{\epsilon}{2}$$

Zatem $\forall_{n > N}: \rho(x_n, g) \leq \rho(x_n, x_M) + \rho(x_M, g) < \frac{\epsilon}{2} + \frac{\epsilon}{2} = \epsilon$, czyli $g \in X$ jest granicą ciągu $(x_n)$, a więc $(x_n)$ jest zbieżny w $(X, \rho)$.
\end{proof}



\subsection{Twierdzenie Baire'a}

\begin{definition}[Zbiór mizerny (zbiór pierwszej kategorii)]
W zupełnej przestrzeni metrycznej zbiorem pierwszej kategorii nazywamy zbiór będący przeliczalną sumą zbiorów nigdziegęstych.
\end{definition}

\begin{definition}[Zbiór rezydualny (zbiór drugiej kategorii)]
\label{definicja_zbioru_rezydualnego}
Zbiór nazywamy rezydualnym, jeżeli zawiera przekrój przeliczalnie wielu gęstych zbiorów otwartych.
\end{definition}



\begin{theorem}[Twierdzenie Baire'a]
\label{twierdzenie_bairea}
W niepustej, zupełnej przestrzeni metrycznej $X$, przekrój przeliczalnie wielu gęstych zbiorów otwartych jest zbiorem gęstym.
\end{theorem}


\begin{lemma}[Przekrój zbiorów rezydualnych jest rezydualny] \label{przekroj_rezydualnych_jest_rezydualny}
W przestrzeni metrycznej zupełnej, przekrój przeliczalnie wielu zbiorów rezydualnych jest rezydualny. W szczególności niepusty.
\end{lemma}



\subsection{Definicje dotyczące układów dynamicznych}

\begin{definition}[Układ dynamiczny]
Układem dynamicznym nazywamy parę $(X, T)$, gdzie $X$ jest zbiorem a $T:X \rightarrow X$ przekształceniem (odwzorowaniem).
\end{definition}

\begin{definition}[Orbita]
Orbita punktu $x$ to ciąg $x, Tx, T^2x, T^3x,\ldots$, czyli 
$$O_T(x) = \{T^nx: n=0,1,2,\ldots\}.$$

Jeśli $T$ jest odwracalne (czyli różnowartościowe i ,,na''), to możemy rozważać orbitę obustronną:
$$\{T^nx: n \in \mathbb{Z}\}.$$
\end{definition}

\begin{definition}[Iterata przekształcenia]
Iterata przekształcenia $T$ to potęga przekształcenia $T^n$, $n \in \mathbb{N}$.
\end{definition}

\begin{definition}[Punkt stały]
Punkt $x$ jest stały, gdy $Tx = x$.
\end{definition}

\begin{definition}[Punkt okresowy odwzorowania]
Punkt $x$ jest okresowy, gdy jest stały dla jakiejś iteraty, tzn. gdy $\exists_{n \in \mathbb{N}} \, T^nx = x$.
\end{definition}


\begin{definition}[Orbita okresowa]
TODO
\end{definition}

\begin{definition}[Niezmienniczość zbioru ze względu na odwzorowanie]
Zbiór $A \subset X$ nazywamy niezmienniczym ze względu na odwzorowanie $f: X \rightarrow X$ jeżeli $f(A) \subset A.$
\end{definition}



\subsection{Entropia topologiczna}

\subsubsection{Definicja Adlera, Konheima i McAndrew}

\begin{definition}[Entropia topologiczna pokrycia]
Na podstawie \cite{misiurewicz1993}.
Rozważamy układ dynamiczny $(X, f)$, gdzie $X$ jest zwartą przestrzenią metryczną. Niech $\mathcal{A}$ będzie otwartym pokryciem $X$.
Wprowadzamy oznaczenia:

$$\bigvee_{i=1}^n \mathcal{A}_i = \{A_1 \cap A_2 \cap \ldots \cap A_n : A_1 \in \mathcal{A}_1, A_2 \in \mathcal{A}_2, \ldots, A_n \in \mathcal{A}_n, A_1 \cap A_2 \cap \ldots \cap A_n \neq \emptyset \},$$ 
gdzie $\mathcal{A}_1, \mathcal{A}_2, \ldots, \mathcal{A}_n$ są otwartymi pokryciami $X$.

$$f^{-n}(\mathcal{A}) = \{f^{-n}(A):A\in \mathcal{A}\},$$

$$\mathcal{A}^n = \bigvee_{i=0}^{n-1}f^{-i}(\mathcal{A}).$$

Ponadto niech $\mathcal{N}(\mathcal{A})$ będzie najmniejszą możliwą mocą podpokrycia wybranego z $\mathcal{A}$ (tzn. licznością podzrodziny rodziny $\mathcal{A}$, która również jest pokryciem $X$). Moc tej rodziny możemy utożsamiać z licznością ponieważ zawsze istnieje podpokrycie skończone, gdyż $X$ jest zwarta.

Wtedy entropią topologiczną odwzorowania $f$ przy pokryciu $\mathcal{A}$ nazywamy wartość:
$$h(f, \mathcal{A}) = \lim_{n \rightarrow \infty}{\frac{1}{n}\log\mathcal{N}(\mathcal{A}^n)}$$
\end{definition}

\begin{definition}[Entropia topologiczna Adler, Konheim, McAndrew]
Na podstawie \cite{misiurewicz1993}.
Adler, Konheim i McAndrew entropię topologiczną odwzorowania $f$ definiują następująco: 
$$h(f) = \sup h(f, \mathcal{A}),$$ gdzie supremum jest brane po wszystkich otwartych pokryciach $\mathcal{A}$ przestrzeni $X$.
\end{definition}


\subsubsection{Definicja Bowena}
\cite{ruette2017chaos}

Niech $(X, \rho)$ będzie zwartą przestrzenią metryczną, a $f: X \rightarrow X$ ciągłym odwzorowaniem (TODO: ma byc uniformly continous - tak jest u Bowena, jest tak ze zwartości $X$, napisać to tutaj?).

\begin{definition}[Kula Bowena]
Niech  $\epsilon > 0$ oraz $n \geq 1$. Kulę Bowena rzędu $n$ o środku w punkcie $x \in X$ i promieniu $\epsilon$ definiujemy następująco

\begin{equation}
\begin{aligned}
B_n(x,\epsilon) & \coloneqq \{y \in X : \rho(f^k(x), f^k(y)) \leq \epsilon, k \in \{0, 1, \ldots, n-1\}\} \\
& = \bigcap_{i=0}^{n-1}f^{-i}(\overline{K}(f^i(x), \epsilon)).
\end{aligned}
\end{equation}

\end{definition} 


\begin{definition}[Zbiór $(n,\epsilon)$-rozdzielony]
Mówimy, że zbiór $E \subset X$ jest $(n, \epsilon)$-rozdzielony jeżeli dla każdych $x, y \in E,$ takich że $x \neq y$ istnieje $k \in \{0, 1, \ldots, n-1\},$ takie że $\rho(f^k(x), f^k(y)) > \epsilon.$

Maksymalną moc zbioru $(n, \epsilon)$-rozdzielonego oznaczamy przez $s_n(f, \epsilon).$
\end{definition}


\begin{definition}[Zbiór $(n, \epsilon)$-rozpinający]
Zbiór $E$ nazywamy $(n,\epsilon)$-rozpinającym jeżeli $X \subset \bigcup_{x \in E}B_n(x, \epsilon)$.

Minimalną moc zbioru $(n, \epsilon)$-rozpinającego oznaczamy przez $r_n(f, \epsilon).$

\end{definition}


\begin{definition}[Entropia topologiczna Bowena]
Niech $(X, f)$ będzie układem dynamicznym.
Wtedy

$$h_{\textrm{top}}(f) = \lim_{\epsilon \rightarrow 0} \limsup_{n \rightarrow +\infty} \frac{1}{n} \log s_n(f, \epsilon) = \lim_{\epsilon \rightarrow 0} \limsup_{n \rightarrow +\infty} \frac{1}{n} \log r_n(f, \epsilon).$$

\end{definition}



\textcolor{red}{TODO: entropia topologiczna - dowod rownowaznosci
 (ksiazka misiurewicz combinatorial dynamics and entropy in dimension one oraz ksiazka ruette rozdzial 4)}
 
 
\textcolor{red}{TODO: nowe} 
\begin{lemma}[Złożenie dwóch funkcji niemalejących jest niemalejące]
\label{zlozenie_dwoch_niemalejacych_jest_niemalejace}
Niech $f: [0, 1] \rightarrow [0, 1]$ i $g: [0, 1] \rightarrow [0,1]$ będą funkcjami niemalejącymi, wówczas ich złożenie $f \circ g$ również jest funkcją niemalejącą. 
\end{lemma}

\begin{proof}
Weźmy dowolne $x, y \in [0,1]$, niech $y < x$. Wówczas, ponieważ $g$ jest niemalejąca zachodzi nierówność $g(y) \leq g(x)$, z faktu, że $f$ również jest niemalejąca otrzymujemy $f(g(y)) \leq f(g(x)).$ Czyli $(f \circ g)(y) \leq (f \circ g)(x).$
\end{proof} 

\begin{corollary}[Złożenie skończenie wielu funkcji niemalejących jest niemalejące]
\label{zlozenie_skonczenie_wielu_niemalejacych_jest_niemalejace}
Wniosek ten otrzymujemy z lematu \ref{zlozenie_dwoch_niemalejacych_jest_niemalejace} przez zastosowanie indukcji matematycznej.
\end{corollary}

\begin{corollary}[Iteraty funkcji niemalejącej są niemalejące]
\label{iteraty_niemalejacych_sa_niemalejace}
Niech $f: [0, 1] \rightarrow [0, 1]$ będzie funkcją niemalejącą, wówczas dla każdego $i \in \mathbb{N}$ funkcje $f^{i}$ są niemalejące.
\end{corollary}
\textcolor{red}{TODO: end nowe} 

 
 
\begin{theorem}[Entropia topologiczna odwzorowań monotonicznych]
\label{entropia_topologiczna_odwzorowan_monotonicznych}

Niech $I = [0, 1]$ i niech $f: I \rightarrow I$ będzie funkcją monotoniczną (niekoniecznie ściśle).
Wówczas entropia topologiczna $f$ jest równa zero.
\end{theorem}

\begin{proof}
Rozważmy dwa przypadki: osobno dla $f: I \rightarrow I$ niemalejącej i nierosnącej.
\item[Przypadek 1: $f: I \rightarrow I$ niemalejąca.]

Skorzystamy z definicji entropii topologicznej Bowena. Pokażemy, że 
$$\lim_{\epsilon \rightarrow 0} \limsup_{n \rightarrow +\infty} \frac{1}{n} \log s_n(f, \epsilon) = 0.$$

% Pokażemy, że $\forall_{n \in \mathbb{N}, \epsilon>0} \, s_n(f, \epsilon) \leq \frac{n}{\epsilon}.$

Ustalmy dowolne $n \in \mathbb{N}$ i $\epsilon > 0.$ Podzielmy odcinek $[0, 1]$ na $m = \lfloor 1/\epsilon \rfloor + 1 \leq \frac{1}{\epsilon} + 1$ odcinków o (niekoniecznie jednakowych) długościach równych $d_{i,\epsilon} < \epsilon$ w taki sposób aby otrzymać rodzinę odcinków: $\{ [0, d_{1,\epsilon}), [d_{1,\epsilon}, d_{1,\epsilon} + d_{2, \epsilon}), \ldots, [\sum_{i=1}^{m-1} d_{i, \epsilon}, 1] \}.$ Oznaczmy krańce tych odcinków przez $a_0=0, \, a_1 = d_{1,\epsilon}, \, a_2 = d_{1,\epsilon} + d_{2,\epsilon}, \, \ldots, \, a_m = 1.$ 

Zdefiniujmy $A_i \coloneqq \{f^{-i}(a_0), f^{-i}(a_1), \ldots, f^{-i}(a_m)\}$, dla $i \in \{0, 1, \ldots, n-1\}$

Niech $B = \bigcup_{i=0}^{n-1} A_i.$ Każdy ze zbiorów $A_i$ zawiera $m+1$ elementów, zatem $|B| \leq n \cdot (m+1)$ Uporządkujmy elementy zbioru $B$ w niemalejący ciąg $(b_j)_{j=0}^{|B|-1}.$ Podzielmy $I$ na odcinki postaci $B_j = [b_{j-1}, b_{j}]$ dla $j \in \{1, 2, \ldots, |B|-1\}.$  

Pokażę, że dowolne dwa punkty $x, y$ należące do tego samego odcinka $B_j$ nie mogą jednocześnie należeć do tego samego zbioru $(n, \epsilon)$-rozdzielonego.

Weźmy dowolny odcinek $B_j = [b_{j-1}, b_{j}] \, j \in \{1, 2, \ldots, |B|-1\}.$ Wybierzmy z niego dwa dowolne punkty $x, y$, $x \neq y$. Bez straty ogólności, niech $x > y.$ Oczywiście nie istnieje punkt $b_k \in B$ dla którego $y < b_k < x.$ Z konstrukcji $B$ widzimy natychmiast, że $\forall_{i \in \{0, 1, \ldots, n-1\}} \, \exists_{j \in \{1,2,\ldots,m\}} \, f^{-i}(a_{j-1}) \leq y < x \leq f^{-i}(a_j)$. Z monotoniczności funkcji $f$ i~wniosku \ref{iteraty_niemalejacych_sa_niemalejace} otrzymujemy 
$$f^i(f^{-i}(a_{j-1})) =  a_{j-1} \leq f^i(y) \leq f^i(x) \leq a_j = f^i(f^{-i}(a_j)).$$
Ponieważ $\rho(a_{j-1}, a_j) = d_{j,\epsilon} < \epsilon$ to  $\rho(f^i(x), f^i(y)) < \epsilon.$ A więc $x$ i $y$ nie należą do jednego zbioru $(n, \epsilon)$-rozdzielonego.



Przypomnijmy, że $s_n(f, \epsilon)$ oznacza maksymalną moc zbioru $(n, \epsilon)$-rozdzielonego. 

% Aby dwa punkty należały do zbioru $(n, \epsilon)$-rozdzielonego, nie mogą one należeć do jednej kuli Bowena o promieniu mniejszym bądź równym $\frac{\epsilon}{2}.$ 

Z wcześniejszych rozważań wynika, że do dowolnego zbioru $(n, \epsilon)$-rozdzielonego może należeć co najwyżej jeden punkt z każdego z odcinków $B_j$, odcinków tych jest $|B|-1.$ Zatem $s_n(f, \epsilon) \leq |B|-1 \leq n \cdot (m+1) - 1 \leq n \cdot (\frac{1}{\epsilon}+1+1) - 1 = \frac{n}{\epsilon} + \frac{n}{2} - 1.$



Ostatecznie otrzymujemy:
\begin{equation}
\begin{aligned}
h_{\textrm{top}}(f) & = \lim_{\epsilon \rightarrow 0} \limsup_{n \rightarrow +\infty} \frac{1}{n} \log s_n(f, \epsilon) \leq \lim_{\epsilon \rightarrow 0} \limsup_{n \rightarrow +\infty} \frac{1}{n} \log ( n \cdot (m+1) - 1 ) \\
& \leq \lim_{\epsilon \rightarrow 0} \limsup_{n \rightarrow +\infty} \frac{1}{n} \log ( \frac{n}{\epsilon} + \frac{n}{2} - 1 ) \leq \lim_{\epsilon \rightarrow 0} \limsup_{n \rightarrow +\infty} \frac{1}{n} \log ( \frac{2 \cdot n + \epsilon \cdot n}{2 \cdot \epsilon}) \\
& = \lim_{\epsilon \rightarrow 0} \limsup_{n \rightarrow +\infty} (\frac{1}{n} \log (n \cdot (2 + \epsilon)) - \frac{1}{n} \log ( 2 \cdot \epsilon )) = \lim_{\epsilon \rightarrow 0} 0 = 0
\end{aligned}
\end{equation}
\textcolor{red}{
Entropia z definicji jest nieujemna, więc $h_{\textrm{top}}(f) = 0.$
}

\item[Przypadek 2: $f: I \rightarrow I$ nierosnąca.]
\textcolor{red}{TODO: nowe!!!}

Jeżeli $f$ jest nierosnąca, to $f^2$ jest niemalejąca. Zatem z przypadku pierwszego $h_{\textrm{top}}(f^2) = 0$. Dla entropii topologicznej zachodzi następująca równość: 
$$h_{\textrm{top}}(f^n) = n \cdot h_{\textrm{top}}(f),$$
zatem 
$0 = h_{\textrm{top}}(f^2) = 2 \cdot h_{\textrm{top}}(f),$
czyli
$h_{\textrm{top}}(f) = 0.$

\end{proof}
 
 
 
 
\begin{theorem}[Dolne ograniczenie na entropię na odcinku] \label{dolne_ograniczenie_entropii_na_odcinku}
\cite{ruette2017chaos}
Jeżeli $f$ jest topologicznie tranzytywne, to entropia topologiczna $h(f) \geq \frac{1}{2} \log(2)$.
TODO
\end{theorem}

\begin{proof}
\textcolor{red}{TODO: jak zabraknie czasu to tylko zacytowac artykul}
\end{proof}



\subsection{Chaos}


Pomimo dużej popularności chaotycznych układów dynamicznych, przez długi czas nie było jednej powszechnie akceptowanej definicji chaosu. W niniejszej pracy prezentujemy dwie spośród powszechnie stosowanych: chaos w sensie Devaneya oraz w sensie Li-Yorke'a. W dalszej części pracy korzystać zajmować się będziemy odwzorowaniami chaotycznymi w sensie Devaneya.

\begin{definition}[Topologiczna tranzytywność]
\cite{onDeveneyDefinitionOfChaos}
Mówimy, że $f$ jest topologicznie tranzytywne w przestrzeni metrycznej $X$, gdy dla każdych niepustych, otwartych podzbiorów $U, V$ przestrzeni $X$ istnieje liczba naturalna $k$ taka, że przekrój $f^k(U) \cap V$ jest niepusty.
\end{definition}

\begin{definition}[Wrażliwość na warunki początkowe]
\cite{balibrea2003topological}
Niech $(X, \rho)$ będzie przestrzenią metryczną, a $f: X \rightarrow X$ odwzorowaniem ciągłym.
Mówimy, że $f$ jest wrażliwe na warunki początkowe jeżeli  istnieje $\delta > 0,$ taki że dla każdego $x \in X$ i dla każdego otoczenia $U$ punktu $x$, istnieją $y \in U$ i $n \geq 0$ dla których zachodzi $\rho(f^n(x), f^n(y)) > \delta$.
\end{definition}

\begin{definition}[Chaos w sensie Devaneya]
Niech $X$ będzie przestrzenią metryczną. Ciągłe odwzorowanie $f: X \rightarrow X$ nazywamy chaotycznym (w sensie Devaneya) na $X$ jeżeli:
\begin{enumerate}
\item $f$ jest tranzytywne,
\item zbiór punktów okresowych $f$ jest gęsty w $X$,
\item $f$ jest wrażliwe na warunki początkowe.
\end{enumerate}
\end{definition}


\begin{theorem}
[Warunki dostateczne chaotyczności w sensie Devaneya]
\label{warunki_dostateczne_chaotycznosci_devaneya}
\cite{onDeveneyDefinitionOfChaos}
Niech $(X, \rho)$ będzie przestrzenią metryczną oraz niech $X$ nie będzie zbiorem skończonym.
Jeżeli $f: X \rightarrow X$ jest tranzytywne i posiada gęsty zbiór punktów okresowych to $f$ jest wrażliwe na warunki początkowe.
\end{theorem}

\begin{proof}
\cite{onDeveneyDefinitionOfChaos}
Załóżmy, że $f: X \rightarrow X$ jest tranzytywne i posiada gęsty zbiór punktów okresowych.

W pierwszej kolejności zauważmy, że istnieje liczba $\delta_0 > 0$ taka, że dla każdego $x \in X$ istnieje punkt okresowy $q \in X$, którego orbita $O(q)$ jest w odległości co najmniej $\frac{\delta_0}{2}$ od $x$: 

Wybierzmy dwa dowolne punkty okresowe $q_1$ i $q_2$ o rozłącznych \textcolor{red}{(TODO: Trzeba uzasadnic dlaczego moge wybrac w punkty o rozlacznych orbitach?, moze napisac tez dlaczego potrzebne jest zalozenie o nieskonczonosci X - pokazac jakis degenerate case)} orbitach $O(q_1), O(q_2)$. Oznaczmy przez $\delta_0$ odległość między $O(q_1)$ a $O(q_2)$. Wówczas z nierówności trójkąta, każdy punkt $x \in X$ jest w odległości co najmniej $\frac{\delta_0}{2}$ od  jednej z dwóch wybranych orbit.
Pokażemy, że $f$ jest wrażliwe na warunki początkowe ze stałą (ang. \textit{sensitivity constant}) $\delta = \frac{\delta_0}{8}.$

Weźmy dowolny punkt $x \in X$ i niech $N$ będzie pewnym otoczeniem punktu $x$. Ponieważ zbiór punktów okresowych odwzorowania $f$ jest gęsty, to istnieje punkt okresowy $p$ należący do przekroju $U = N \cap K(x, \delta)$. Niech $n$ oznacza okres punktu $p$. Z wcześniejszych rozważań wiemy, że istnieje punkt okresowy $q \in X$, którego orbita $O(q)$ jest w odległości co najmniej $4\delta$ od $x$. Oznaczmy

$$V = \bigcap_{i=0}^n f^{-i}(K(f^i(q), \delta)).$$

Zbiór $V$ jest otwarty i niepusty, gdyż $q \in V.$ W związku z tym, ponieważ $f$ jest tranzytywne, to istnieją $y \in U$ i $k \in \mathbb{N},$ takie że $f^k(y) \in V$.

Niech teraz $j$ będzie częścią całkowitą z $\frac{k}{n} + 1$. Wtedy $1 \leq nj - k \leq n$. Z konstrukcji mamy

$$f^{nj}(y) = f^{nj-k}(f^k(y)) \in f^{nj-k}(V) \subseteq K(f^{nj-k}(q), \delta).$$

Liczba $n$ jest okresem $p$, więc $f^{nj}(p) = p$, a z nierówności trójkąta:

$$\rho(x, f^{nj-k}(q)) \leq \rho(x, p) + \rho(p, f^{nj}(y)) + \rho(f^{nj}(y), f^{nj-k}(q)),$$

po przekształceniach dostajemy:

\begin{equation}
\begin{aligned}
\rho(f^{nj}(p), f^{nj}(y)) & = \rho(p, f^{nj}(y)) \\
& \geq \rho(x, f^{nj-k}(q)) - \rho(f^{nj-k}(q), f^{nj}(y)) - \rho(p, x),
\end{aligned}
\end{equation}
Następnie, ponieważ $p \in K(x, \delta)$ oraz $f^{nj}(y) \in K(f^{nj-k}(q), \delta)$ otrzymujemy

$$\rho(f^{nj}(p), f^{nj}(y) > 4\delta - \delta - \delta = 2\delta.$$

Ostatecznie, korzystając z nierówności trójkąta dostajemy, że albo $\rho(f^{nj}(x), f^{nj}(y)) > \delta$, albo $\rho(f^{nj}(x), f^{nj}(p)) > \delta.$ W obu przypadkach znaleźliśmy punkt, którego $nj$-ta iterata jest w odległości większej od $\delta$ od $f^{nj}(x).$
\end{proof}



\begin{definition}[Chaos w sensie Li-Yorke'a]
\cite[s.~25]{aulbach2001three}
Ciągłe odwzorowanie $f: X \rightarrow X$ na zwartej przestrzeni metrycznej $(X, \rho)$ jest chaotyczne (w sensie Li-Yorke'a) jeżeli istnieje nieprzeliczalny podzbiór $S$ przestrzeni $X$ o następujących własnościach:
\begin{enumerate}
\item $\limsup_{n \rightarrow \infty} \rho(f^n(x), f^n(y)) > 0$ dla każdych $x,y \in S, \, x \neq y,$
\item $\liminf_{n \rightarrow \infty} \rho(f^n(x), f^n(y)) = 0$ dla każdych $x,y \in S, \, x \neq y,$
\item $\limsup_{n \rightarrow \infty} \rho(f^n(x), f^n(p)) > 0$ dla każdych $x \in S, \, p \in X, \, p \textrm{ - okresowy.}$
\end{enumerate}
\end{definition}



\subsection{Odwzorowania trójkątne}


\begin{definition}[Odwzorowanie trójkątne]
Rozważmy ,,prostokąt'' $X \times Y$, gdzie $X, Y$ są przestrzeniami metrycznymi. 
Odwzorowaniem trójkątnym nazywamy funkcję $F: X \times Y \rightarrow X \times Y$ postaci:
$$F(x,y) = (f(x), g(x,y)).$$
\end{definition}


\begin{definition}[Ciągłość odwzorowań trójkątnych]
W pracy rozważać będziemy wyłącznie odwzorowania ciągłe, tj. takie trójkątne $F$ dla których $f \in C(X)$ i $g$ jest ciągłą funkcją z $X \times Y$ w $Y$. 
Zbiór takich odwzorowań oznaczać będziemy $C_\triangle(X \times Y)$
Zamiast $g(x, y)$ będziemy pisać $g_x(y)$, gdzie $g_x: Y \rightarrow Y$ jest rodziną ciągłych przekształceń zależną w sposób ciągły od $x \in X$.
\end{definition}


\begin{definition}[Odwzorowanie bazowe]
Odwzorowanie $f$ nazywamy odwzorowaniem bazowym (ang. \textit{basis map}) odwzorowania $F = (f, g_x)$
\end{definition}

\begin{definition}[Odwzorowania włóknowe]
Odwzorowania $g_x$ dla $x \in X$ nazywamy odwzorowaniami włóknowymi (ang. \textit{fibre maps}) odwzorowania $F = (f, g_x)$.
\end{definition}


\begin{theorem}[O równości entropii odwzorowania trójkątnego i jego odwzorowania bazowego]
\label{rownosc_entropii_gdy_wlokna_monotoniczne}
Niech $X$ będzie zwartą przestrzenią metryczną, $I = [0, 1]$, $F = (f, g_x) \in C_\triangle(X \times I).$ Wtedy, jeżeli $\forall_{x \in X}$ funkcje włóknowe $g_x$ są (niekoniecznie ściśle) monotoniczne, to entropia topologiczna odwzorowania $F$ jest równa entropii topologicznej jego odwzorowania bazowego $f$, czyli:

$$h(F) = h(f)$$
\end{theorem}

\begin{proof}
%\cite{do_dowodu_rownosci_topologicznej_entropii_2_twierdzenieD}

Rozważamy $X \times I$, oraz $X.$ Są to zwarte przestrzenie metryczne (TODO: pokazać zwartość? Jaką metrykę mamy na $X \times I$? max($\rho_x, \rho_y$)?).
Zdefiniujmy odwzorowanie $\pi: X \times I \rightarrow X$, $\pi(x,a) = x$. Oczywiście $\pi$ jest ciągłe i ,,na''.
Wówczas $F: X \times I \rightarrow X \times Y$, $f: X \rightarrow X$ oraz $\pi: X \times I \rightarrow X$ (,,na'') są ciągłymi odwzorowaniami. 

Weźmy dowolny punkt $(x, a) \in X \times I$:
$$(\pi \circ F)(x, a) = \pi(F(x,a)) = \pi(f(x), g_x(a)) = f(x),$$
$$(f \circ \pi)(x, a) = f(\pi(x,a)) = f(x),$$
a więc $\pi \circ F = f \circ \pi.$


Na podstawie twierdzenia 17 z pracy  \cite[s.~103]{bowen1971entropy__do_dowodu_rownosci_topologicznej_entropii_2_twierdzenie17} otrzymujemy, że 

$$h(F) \leq h(f) + \sup_{x \in X} h(F, \pi^{-1}(x))$$

\textcolor{red}{Ustalmy $x \in X$, $h(F, \pi^{-1}(x)) = h(F, \{x\} \times I) =???= h(g_x)$ TODO: nie do konca rozumiem zapis $h(F, A)$, gdzie $A$ jest zbiorem, w pracy Bowena zbiór zwarty $K$ występuje przy definicji 2 na stronie 402 i nie jest powiedziane wprost, że ograniczamy się do tego zbioru, tylko, że przy zbiorach $(\epsilon, n)$-rozdzielonych ograniczamy się do największej mocy zbioru $(\epsilon, n)$-rozdzielonego zawartego w K. Chyba powinienem wykazać, że to sprowadza się do entropii na zbiorze K (w naszym przypadku do włókna z odwzorowaniem $g_x$)}

Z twierdzenia \ref{entropia_topologiczna_odwzorowan_monotonicznych} i faktu, że wszystkie odwzorowania włóknowe $g_x$ są monotoniczne mamy :
$$\forall_{x \in X} \, h(F, \pi^{-1}(x)) = h(g_x) = 0,$$
czyli również:
$$\sup_{x \in X} h(F, \pi^{-1}(x)) = 0.$$

Ostatecznie dostajemy:

$$h(F) \leq h(f)$$

\textcolor{red}{TODO: pokazac ze $h(F) \geq h(f)$}

\end{proof}



\subsection{Definicje utworzone na potrzeby dowodu twierdzenia o rozszerzaniu}
Definicje i własności odwzorowań trójkątnych podajemy na podstawie pracy \cite{balibrea2003topological}.

Na potrzeby dowodu wprowadźmy pojęcia odległości między odwzorowaniami oraz dwie funkcje: $\textrm{pr}_1(x, y)$ i $\textrm{pr}_2(x, y)$. 

\begin{definition}[Metryka na przestrzeni ciągłych przekształceń zwartej przestrzeni metrycznej]
Niech $(M, \sigma)$ będzie zwartą przestrzenią metryczną, rozważmy odwzorowania $h,k \in C(M)$. Odległość między nimi zdefiniujmy jako $\max_{m \in M} \sigma(h(m), k(m))$ i oznaczmy ją jako $d_1(h,k)$.
\end{definition}


\begin{definition}[Metryka na przestrzeni odwzorowań trójkątnych]
Odległość między odwzorowaniami trójkątnymi definiujemy wówczas następująco: Niech $(X, \rho)$ i $(Y, \tau)$ będą zwartymi przestrzeniami metrycznymi a $F(x,y) = (f(x), g_x(y))$ i $\Phi(x,y) = (\phi(x), \psi_x(y))$ trójkątnymi odwzorowaniami należącymi do $C_\triangle(X \times Y)$. Odległość definiujemy wówczas jako 
\begin{align*}
d_2(F, \Phi) = \max_{(x,y) \in X \times Y} \max\{\rho(f(x),\phi(x)), \tau(g_x(y), \psi_x(y))\} \\ 
= \max\{d_1(f,\phi), \max_{x \in X}d_1(g_x, \psi_x)\}
\end{align*}
Zauważmy, że jak wynika z lematów \ref{przestrzen_ciaglych_jest_zupelna} i \ref{ciagle_trojkatne_tworza_przestrzen_metryczna_zupelna} przestrzenie metryczne $(C(X), d_1)$ oraz $(C_\triangle(X \times Y), d_2)$ są zupełne i odpowiednie topologie na nich są topologiami zbieżności jednostajnej.
\end{definition}


\begin{lemma}[Przestrzeń $(C(X), d_1)$ jest zupełna]\label{przestrzen_ciaglych_jest_zupelna}
Niech $(X, \rho)$ będzie zwartą przestrzenią metryczną. Przestrzeń metryczna $(C(X), d_1)$ jest zupełna i topologia na niej jest topologią zbieżności jednostajnej.
\end{lemma}

\begin{proof}
{\color{red} !!!NOWE!!!

Musimy pokazać, że każdy ciąg Cauchy'ego w przestrzeni $(C(X), d_1)$ jest zbieżny do elementu tej przestrzeni.
Niech $(f_n)_{n=0}^{\infty} \subset C(X),$ będzie ciągiem Cauchy'ego, czyli 
\begin{equation}
\label{warunek_cauchyego_na_przestrzeni_funkcji_ciaglych}
\forall_{\epsilon > 0} \, \exists_{N \in \mathbb{N}} \, \forall_{n,m > N} \, d_1(f_n, f_m) = \max_{x \in X} \rho\big(f_n(x), f_m(x)\big) < \epsilon.
\end{equation}

Weźmy dowolny $x_0 \in X$, oczywiście 
\begin{equation}
\label{metryka_w_punkcie_mniej_rowna_niz_metryka_zbieznosci_jednostajnej}
\forall_{n,m \in \mathbb{N}} \, \rho\big(f_n(x_0), f_m(x_0)\big) \leq \max_{x \in X} \rho\big(f_n(x), f_m(x)\big).
\end{equation}

Podstawiając \ref{metryka_w_punkcie_mniej_rowna_niz_metryka_zbieznosci_jednostajnej} do \ref{warunek_cauchyego_na_przestrzeni_funkcji_ciaglych} otrzymujemy, że dla każdego $x \in X$ ciąg $\big(f_n(x)\big)_{n=0}^{\infty}$ jest ciągiem Cauchy'ego w przestrzeni $(X, \rho).$ Ponieważ $(X, \rho)$ jest przestrzenią zwartą a więc zupełną (twierdzenie \ref{zwarta_jest_zupelna}), to dla każdego $x \in X$ ciąg $\big(f_n(x)\big)_{n=0}^{\infty}$ jest zbieżny.

Oznaczmy teraz $f(x) = \lim_{n \rightarrow \infty} f_n(x).$ 


Pozostało pokazać, że ciąg $(f_n)_{n=0}^{\infty} \subset C(X)$ zbiega do $f$ w przestrzeni $(C(X), d_1).$ Weźmy dowolny $\epsilon > 0.$ Pokażemy, że istnieje $n_0 \in \mathbb{N},$ takie że dla każdego $n > n_0$ zachodzi

$$d_1(f_n, f) = \max_{x \in X} \rho\big(f_n(x), f(x)\big) < \epsilon$$

Z definicji odwzorowania $f$, dla każdego $x \in X$ ciąg $\big(f_n(x)\big)_{n=0}^{\infty}$ jest zbieżny do $f(x)$, a więc istnieje $n_x \in \mathbb{N},$ takie że $\forall_{n > n_x} \, \rho\big(f_n(x), f(x)\big) < \epsilon.$ Wystarczy więc wziąć $n_0 = \max_{x \in X} n_x.$

Zauważmy, że $d_1$ jest metryką zbieżności jednostajnej. Odwzorowanie $f$ jest granicą ciągu odwzorowań ciągłych w metryce zbieżności jednostajnej, a więc $f \in C(X).$
}
\end{proof}

\begin{lemma}[Ciągłe odwzorowania trójkątne tworzą przestrzeń metryczną zupełną]
\label{ciagle_trojkatne_tworza_przestrzen_metryczna_zupelna}
Przestrzeń metryczna $(C_\triangle(X \times Y), d_2)$ jest zupełna i topologia na niej jest topologią zbieżności jednostajnej.
\end{lemma}

\begin{proof}
TODO
\end{proof}


\begin{definition}[$\textrm{pr}_1(x, y)$, $\textrm{pr}_2(x, y)$]
Dla $(x, y) \in X \times Y$ niech $\textrm{pr}_1(x, y) = x$ i $\textrm{pr}_2(x, y) = y$. Odwzorowanie identycznościowe na $Y$ będziemy oznaczać przez $\textrm{Id}_Y$ lub krótko $\textrm{Id}$. W dalszej części pracy przestrzeń $Y$ będzie odcinkiem rzeczywistym $I = [0, 1]$.
\end{definition}





\section{Lematy}

Na potrzeby kolejnych lematów wprowadźmy następujące definicje i oznaczenia.

Niech $(X, \rho)$ będzie zwartą przestrzenią metryczną bez punktów izolowanych oraz niech $f \in C(X)$ będzie odwzorowaniem chaotycznym w sensie Devaneya.
Przez $I$ oznaczmy odcinek domknięty $[0, 1].$
Niech $\mathcal{F}$ będzie zbiorem wszystkich odwzorowań $F = (f, g_x)$ ze zbioru $C_\triangle(X \times I)$ spełniających następujące warunki:
\begin{enumerate}
\item Odwzorowanie bazowe $f$ spełnia założenia twierdzenia \ref{twierdzenie_glowne}.
\item $\forall_{x \in X}$ odwzorowanie $g_x$ jest niemalejące i krańce przedziału I pozostawia niezmienione.
\item $\forall_{x \in P_0}$ $g_x$ jest identycznością 
\end{enumerate}


\begin{lemma} \label{F_jest_niepustym_domknietym_podzbiorem_trojkatnych}
Zbiór $\mathcal{F}$ jest niepustym, domkniętym podzbiorem przestrzeni $C_{\triangle}(X \times I)$, 
\end{lemma}

\begin{proof}
\textcolor{red}{TODO}
\end{proof}


\begin{theorem} \label{tranzytywne_rezydualne_w_F}
Zbiór odwzorowań tranzytywnych jest rezydualny w $\mathcal{F}$.
\cite{alseda1999entropy} - dowod twierdzenia 1.5
\end{theorem}

\begin{proof}
\textcolor{red}{TODO: dowód jest w pracy {alseda1999entropy} - dowod twierdzenia 1.5}
\end{proof}



\begin{lemma} \label{Fso_sa_otwartymi_podzbiorami_F}
Niech $U$ będzie dowolnym zbiorem otwartym zawartym w $X \times I$. Zdefiniujmy zbiór $\mathcal{F}_{SO}$ (,,S'' - stabilny, ,,O'' - okresowy) następująco. Odwzorowanie $G$ należy do $\mathcal{F}_{SO}$ wtedy i tylko wtedy, gdy należy do $\mathcal{F}$, posiada punkt okresowy w $U$ oraz wszystkie dostatecznie bliskie $G$ odwzorowania z $\mathcal{F}$ również posiadają punkt okresowy w $U$ (być może różne od punktów okresowych odwzorowania $G$).
Wówczas zbiór $\mathcal{F}_{SO}$ jest otwartym podzbiorem $\mathcal{F}$.
\end{lemma}

\begin{proof}
\textcolor{red}{TODO}
\end{proof}



\begin{lemma}
\label{w_otoczeniu_niezmienniczego_jest_otwarty_zawarty_w_iteracjach}

Niech $(X, \rho)$ będzie zwartą przestrzenią metryczną bez punktów izolowanych, $f \in C(X)$ odwzorowaniem chaotycznym w sensie Devaneya, $P$ skończonym zbiorem niezmienniczym ze względu na $f$, a $V$ otwartym otoczeniem zbioru $P$.
Wówczas dla każdego $N \in \mathbb{N}$ istnieje niepusty zbiór otwarty $W \subseteq V,$ taki że $W \cup f(W) \cup \ldots \cup f^N(W) \subseteq V$.
\end{lemma}

\begin{proof}
{\color{red} !!!NOWE!!! Pytanie: czy zawsze zlozenie funkcji ciaglych jest ciagle? Korzystam tutaj z tego. Poza tym wygląda na to że wystarczyły jedynie założenia o ciągłości f i skończoności niezmienniczego zbioru P, czy ja czegoś nie przeoczyłem?
W pierwszej kolejności pokażemy prawdziwość tezy dla $N=2.$
$P$ jest zbiorem skończonym, więc możemy wypisać wszystkie jego elementy: 
$$P = \{p_1, p_2, \ldots, p_k\}$$
Ponieważ $V$ jest otwartym otoczeniem zbioru $P$ to istnieje $d>0,$ takie że dla każdego ${0<\epsilon<d}$ zbiór $\bigcup_{i=0}^k K(p_i, \epsilon)$ zawiera się w $V.$ Ustalmy $0<\epsilon<d$ i oznaczmy $U=\bigcup_{i=0}^k K(p_i, \epsilon),$ oczywiście $U \subseteq V.$ Rozważmy dowolny, ustalony punkt $p_i$ ze zbioru $P.$ Wiemy, że $P$ jest niezmienniczy ze względu na odwzorowanie $f$, a więc istnieją punkty $p_{j_1}, p_{j_2}, \ldots, p_{j_N} \in P$ dla których zachodzi $f(p_i) = p_{j_1}, \, f^2(p_i) = p_{j_2}, \, \ldots, \, f^N(p_i) = p_{j_N}.$ Z kolei z ciągłości odwzorowania $f$, również jego iteraty są ciągłe, a z tego wynika, że znajdziemy taką $\delta_i,$ że $\epsilon >\delta_i > 0$ i dla której $f^n(K(p_i, \delta_i)) \subset K(p_{j_n}, \epsilon) \subset U,$ dla każdego $n \in \{1,2,\ldots,N\}.$ Z dowolności wyboru $p_i$ otrzymujemy, że każdy zbiór postaci $f^n(K(p_i, \delta_i)),$ (gdzie $n \in \{1,2,\ldots,N\}$ a $p_i \in P$) zawiera się w $U$. Oznaczmy $W = \bigcup_{i=1}^k K(p_i, \delta_i),$ oczywiście $W$ jest niepustym zbiorem otwartym, który zawiera się w $U$ ponieważ $\delta_i < \epsilon$ dla każdego $i \in \{1,2,\ldots,k\}.$ Ponadto
$$\forall_{n \in \{1,2,\ldots,k\}} \, f^n(W) = f^n(\bigcup_{i=1}^k K(p_i, \delta_i)) = \bigcup_{i=1}^k f^n(K(p_i, \delta_i)) \subset U.$$
Ostatecznie otrzymujemy $W \cup f(W) \cup \ldots \cup f^N(W) \subset U \subseteq V.$

}
\end{proof}



\begin{lemma}[Istnienie orbity okresowej odwiedzającej rodzinę zbiorów otwartych]
\label{lemat_3_glownego_artykulu_istnieje_orbita_okresowa_krojaca_sie_z_rodzina_otwartych}

\textcolor{red}{TODO: przeredagowałem twierdzenie i dowód korzystjąc z pracy \cite[s.~7]{someAspectsOfTopologicalTransitivity} - twierdzenie 2.2.2.16}

Na podstawie \cite[s.~231-232 Lemma 3]{balibrea2003topological} i \cite[s.~7]{someAspectsOfTopologicalTransitivity}
Niech $(X, \rho)$ będzie zwartą przestrzenią metryczną.
Niech $f \in C(X)$ będzie topologicznie tranzytywnym odwzorowaniem, którego zbiór punktów okresowych jest gęsty w $X$. 

Wtedy dla każdej rodziny niepustych zbiorów otwartych $U_1, U_2, \ldots, U_n \subset X$ istnieje orbita okresowa $f$, której przekrój z każdym ze zbiorów $U_i, i=1,2,\ldots,n$ jest niepusty
\end{lemma}

\begin{proof}
Ponieważ odwzorowanie $f$ jest tranzytywne a przestrzeń $X$ jest zwarta, to istnieje punkt $x_0$, którego orbita jest gęsta w $X$. Zatem orbita punktu $x_0$ odwiedza każdy ze zbiorów otwartych $U_i, \, i = 1,2,\ldots,n.$ Teraz korzystając z ciągłości odwzorowania $f$ wystarczy wziąć punkt okresowy $f$ dostatecznie bliski $x_0.$
\end{proof}

\begin{proof}
Alternatywa.
x jest zwartaZwarta => osrodkowa i zupelna + twierdzenie bairea
TODO
Przestrzeń $X$ jest zwarta, a zatem zupełna i ośrodkowa.
\end{proof}



\begin{lemma}
\label{lemat_4_glownego_artykulu}
\cite{alseda1999entropy}
Niech $(X, \rho)$ będzie zwartą przestrzenią metryczną i niech $F = (f, g_x)$ będzie odwzorowaniem należącym do $C_\triangle(X \times I)$ którego wszystkie odwzorowania włóknowe są niemalejące i pozostawiają krańce $I$ niezmienione. Niech $\{a_1, a_2, \ldots, a_n\}$ będzie podzbiorem $X$ oraz dla $i = 1, 2, \ldots, n$ niech $U_i$ będą parami rozłącznymi zbiorami otwartymi, takimi że $a_i \in U_i$. Załóżmy, że $h_i$ są niemalejącymi odwzorowaniami z $C(I)$ pozostawiającymi krańce $I$ niezmienione i spełniającymi $d_1(h_i, g_{a_i}) < \epsilon$ dla pewnego dodatniego $\epsilon$ i każdego $i = 1, 2, \ldots, n$. Wówczas istnieje odwzorowanie $\widetilde{F} = (f, \widetilde{g}_x) \in C_\triangle(X \times I)$ spełniające cztery następujące warunki:

\begin{enumerate}
\item  wszystkie odwzorowania włóknowe $\widetilde{F}$ są niemalejące i pozostawiające krańce $I$ niezmienione,
\item $d_2(F, \widetilde{F}) < \epsilon$,
\item $\widetilde{g}_{a_i} = h_i$ dla $i = 1,2,\ldots,n$,
\item $\widetilde{g}_x = g_x$ dla $x \in X \setminus \bigcup_{i=1}^n U_i$.
\end{enumerate}

\begin{proof}
\cite{alseda1999entropy}
Dla każdego $i = 1, 2, \ldots, n$ niech $V_i \subset U_i$ będzie otwartym otoczeniem $a_i,$ takim że dla pewnego dodatniego $\widetilde{\epsilon} < \epsilon$, $d_1(h_i, g_x) < \widetilde{\epsilon}$ zawsze wtedy gdy $x \in V_i$.

Oznaczmy $U = \bigcup_{i=1}^n U_i$, $V = \bigcup_{i=1}^n V_i$. Niech $u : X \longrightarrow [0,1]$ \textcolor{red}{TODO: cwiczenie - dlaczego funkcja u istnieje?} będzie ciągłą funkcją, przyjmującą wartość $1$ na zbiorze $\{a_1, a_2, \ldots, a_n\}$, natomiast $0$ poza zbiorem $V$. Zastąpmy każde odwzorowanie włóknowe $g_x$ przez $\widetilde{g}_x$, gdzie

\begin{equation} \label{lemat4_rownanie_1}
    \widetilde{g}_x(y) =
    \begin{cases}
        g_x(y) & \text{dla $x \in X \setminus V$,}\\
        g_x(y)(1-u(x))+h_i(y)u(x) & \text{dla $x \in V_i : i \in \{1,2,\ldots,n\}$.}
    \end{cases}
\end{equation}

\textcolor{red}{TODO: Przemyslec, w kontekscie powyzszego zdanie: kombinacja wypukla odwzorowan niemalejacych jest niemalejaca - do czego to tutaj jest porzebne}

dla każdego $y \in I$. Zauważmy ponadto, że dla $x \in V_i$ oraz $i \in \{1,2,\ldots,n\}$ możemy równoważnie napisać

\begin{equation} \label{lemat4_rownanie_2}
\widetilde{g}_x(y) = u(x)(h_i(y) - g_x(y)) + g_x(y).
\end{equation}

Rozważmy odwzorowanie $\widetilde{F} = (f, \widetilde{g}_x)$. Należy ono do $C_\triangle(X \times I)$. Z równości \ref{lemat4_rownanie_1} widzimy, że wszystkie odwzorowania włóknowe $\widetilde{g}_x$ są niemalejące i krańce przedziału $I$ są ich punktami stałymi. Ponadto $\widetilde{g}_{a_i} = h_i$ dla każdego $i$ oraz $\widetilde{g}_x = g_x$ dla $x \in X \setminus V \supset X \setminus U$. Ponieważ dla $x \in V_i$, gdzie $i \in \{1,2,\ldots,n\}$ mamy $d_1(h_i, g_x) < \widetilde{\epsilon}$ oraz $u(x) \in [0,1]$, zatem z równości \ref{lemat4_rownanie_2} otrzymujemy $d_1(g_x, \widetilde{g}_x) < \widetilde{\epsilon}$ dla każdego $x \in V$. Wynika z tego, że $d_2(F, \widetilde(F)) \leq \widetilde{\epsilon} < \epsilon$ (TODO uzasadnic tutaj wewnatrz to ostatnie przejscie od d1 do d2 z definicji d1 i d2), co kończy dowód.

\end{proof}

\end{lemma}



\chapter{Twierdzenie o rozszerzaniu odwzorowań chaotycznych w sensie Devaneya}
\begin{theorem}[O rozszerzaniu]\label{twierdzenie_glowne}
Twierdzenie wraz z dowodem przytaczamy za pracą \cite{balibrea2003topological}
Niech $(X, \rho)$ będzie zwartą przestrzenią metryczną bez punktów izolowanych oraz niech $f \in C(X)$ będzie odwzorowaniem chaotycznym w sensie Devaneya. Wówczas odwzorowanie $f$ można rozszerzyć do odwzorowania $F \in C_{\triangle}(X \times I)$ (to znaczy tak, że $f$ jest odwzorowaniem bazowym dla $F$) w taki sposób, że:
\begin{enumerate}[label=(\roman*)]
\item\label{twierdzenie_glowne_a} $F$ jest również chaotyczne w sensie Devaneya, 
\item\label{twierdzenie_glowne_b} $F$ ma taką samą entropię topologiczną jak $f$, 
\item\label{twierdzenie_glowne_c} zbiory $X \times \{0\}$ i $X \times \{1\}$ są niezmiennicze ze względu na $F$.
\end{enumerate}
\end{theorem}


\begin{proof}[Dowód twierdzenia o rozszerzaniu]
Odwzorowanie $f$ jest chaotyczne w sensie Devaneya, zatem ma gęsty zbiór punktów okresowych, w szczególności istnieje orbita okresowa.
Możemy zatem ustalić okresową orbitę $P_0$ odwzorowania $f$. Ponieważ $P_0$ jest zbiorem skończonym a $X$ nie ma punktów izolowanych, to $P_0$ jest nigdziegęstym domkniętym podzbiorem $X$.

Rozważmy zbiór $\mathcal{F}$ wszystkich odwzorowań $F = (f, g_x)$ ze zbioru $C_\triangle(X \times I)$ spełniających następujące warunki:
\begin{enumerate}
\item\label{dowod_glowny_a} Odwzorowanie bazowe $f$ jest odwzorowaniem z założenia twierdzenia \ref{twierdzenie_glowne}.
\item\label{dowod_glowny_b} $\forall_{x \in X}$ odwzorowanie $g_x$ jest niemalejące i krańce przedziału I pozostawia niezmienione.
\item\label{dowod_glowny_c} $\forall_{x \in P_0}$ $g_x$ jest identycznością 
\end{enumerate}

Warunek \ref{dowod_glowny_b} implikuje, że dla każdego odwzorowania z $\mathcal{F}$ zbiory $X \times \{0\}$ i $X \times \{1\}$ są niezmiennicze, czyli $\forall_{F \in \mathcal{F}}$ zachodzi warunek \ref{twierdzenie_glowne_c}  twierdzenia \ref{twierdzenie_glowne}. Zachodzenie warunku \ref{twierdzenie_glowne_b} twierdzenia \ref{twierdzenie_glowne} dla każdego odwzorowania $F \in \mathcal{F}$ wynika z faktu, że każde odwzorowanie z $C_\triangle(X \times I)$, gdzie $X$ jest zwartą przestrzenią metryczną (a więc w szczególności każde odwzorowanie z $\mathcal{F}$), którego wszystkie odwzorowania włóknowe są niemalejące, ma taką samą entropię topologiczną jak jego odwzorowanie bazowe (twierdzenie \ref{rownosc_entropii_gdy_wlokna_monotoniczne}).
Pozostaje zatem wykazać prawdziwość warunku \ref{twierdzenie_glowne_a}, czyli chaotyczność w sensie Devaneya jakiegoś odwzorowania $F \in \mathcal{F}$. Takie odwzorowanie będzie bowiem łącznie spełniało wszystkie 3 warunki, czyli tezę twierdzenia.

Z lematu \ref{warunki_dostateczne_chaotycznosci_devaneya} wynika, że aby odwzorowanie $F$ było chaotyczne w sense Devaneya potrzeba i wystarcza, żeby spełniało dwa poniższe warunki:
\begin{enumerate}
\item \label{devaney_pierwsza_wlasnosc} $F$ jest topologicznie tranzytywne
\item \label{devaney_druga_wlasnosc} zbiór punktów okresowych odwzorowania $F$ jest gęsty w $(X \times I)$
\end{enumerate}
Jest tak gdyż przestrzeń $(X \times I)$ spełnia założenia lematu, tj. jest przestrzenią nieskończoną i zwartą.

Chcemy wykazać, że istnieje $F \in \mathcal{F}$ będące jednocześnie topologicznie tranzytywne i posiadające gęsty w $(X \times I)$ zbiór punktów okresowych. Z lematu \ref{F_jest_niepustym_domknietym_podzbiorem_trojkatnych} wiemy, że $\mathcal{F}$ jest niepustym, domkniętym podzbiorem przestrzeni $C_{\triangle}(X \times I)$, która jak wynika z lematu \ref{ciagle_trojkatne_tworza_przestrzen_metryczna_zupelna} jest przestrzenią metryczną zupełną. 

Przekrój dwóch zbiorów rezydualnych jest rezydualny (patrz lemat \ref{przekroj_rezydualnych_jest_rezydualny}) a więc niepusty. Wystarczy zatem pokazać, że oba zbiory:
\begin{itemize}
\item zbiór odwzorowań topologicznie tranzytywnych
\item zbiór odwzorowań, których zbiór punktów okresowych jest gęsty w $(X \times I)$
\end{itemize}
są rezydualne w $\mathcal{F}$. Wówczas każde odwzorowanie należące do ich przekroju będzie spełniało wszystkie trzy warunki tezy twierdzenia \ref{twierdzenie_glowne} 

Zbiór odwzorowań tranzytywnych jest rezydualny w $\mathcal{F}$, zostało to udowodnione w twierdzeniu \ref{tranzytywne_rezydualne_w_F}

Pozostało wykazać, że również zbiór odwzorowań posiadających gęsty zbiór punktów okresowych jest rezydualny w $\mathcal{F}$.
Oznaczmy zbiór takich odwzorowań (jednocześnie należących do $\mathcal{F}$) przez $\mathcal{F}_{DP}$.

Niech $\{U_i^X\}_{i=1}^{\infty}$ będzie bazą topologii $X$ i niech $\{U_i^I\}_{i=1}^{\infty}$ będzie zbiorem wszystkich odcinków otwartych o końcach wymiernych, należących do odcinka otwartego $(0, 1)$. Niech $\{U_i\}_{i=1}^{\infty}$ będzie ponumerowaniem zbioru $\{U_i^X \times U_j^I : i,j \in \mathbb{N}\}$. Wtedy każda kula otwarta w $X \times I$ zawiera jakiś spośród otwartych zbiorów $U_i$.

Dla każdego $i=1,2,...$ niech zbiór $\mathcal{F}_{SO}^i$ (,,S'' - stabilny, ,,O'' - okresowy) będzie zdefiniowany następująco. Odwzorowanie G należy do $\mathcal{F}_{SO}^i$ wtedy i tylko wtedy, gdy należy do $\mathcal{F}$, posiada punkt okresowy w $U_i$ oraz wszystkie dostatecznie bliskie $G$ odwzorowania z $\mathcal{F}$ również posiadają punkt okresowy w $U_i$ (być może różne od punktów okresowych odwzorowania $G$). Zbiory $\mathcal{F}_{SO}^i$ są otwartymi podzbiorami $\mathcal{F}$ (patrz lemat \ref{Fso_sa_otwartymi_podzbiorami_F}). Ponieważ $\mathcal{F}_{DP} \supseteq \bigcap_{i=1}^{\infty} \mathcal{F}_{SO}^i$ aby pokazać, że $\mathcal{F}_{DP}$ jest rezydualny w $\mathcal{F}$ wystarczy pokazać, że $\forall_{i \in \mathbb{N}} \, \mathcal{F}_{SO}^i$ jest gęsty w $\mathcal{F}$. (Wynika to z twierdzenia Baire'a \ref{twierdzenie_bairea} oraz definicji \ref{definicja_zbioru_rezydualnego}).

Aby wykazać że każdy zbiór $\mathcal{F}_{SO}^i$ jest gęsty w $\mathcal{F}$ ustalmy dowolne: $i \in \mathbb{N}$, $F=(f,g_x) \in \mathcal{F}$ i $\epsilon > 0$. Pokażemy, że istnieje odwzorowanie $G \in \mathcal{F}_{SO}^i$, którego odległość od $F$ nie przekracza $\epsilon$. Dla uproszczenia sytuacji załóżmy, że $\rho(\textrm{pr}_1(U_i), P_0) > 0$ (Jeżeli tak nie jest, zawsze możemy wziąć zamiast $U_i$ mniejszy prostokąt $U_i^* \subset U_i$).

Weźmy dodatnią liczbę naturalną $N \geq \frac{4}{\epsilon}$. Następnie rozważmy otwarte otoczenie $V$ orbity $P_0$ w przestrzeni $(X, \rho),$ takie że $\rho(\textrm{pr}_1(U_i), V) \geq 0$ oraz $d_1(g_x, \textrm{Id}) < \frac{\epsilon}{4}$ dla każdego $x \in V$ (pamiętamy, że $g_x = \textrm{Id}$ dla $x \in P_0$).

$P_0$ jest skończonym zbiorem niezmienniczym ze względu na $f$, więc na mocy lematu \ref{w_otoczeniu_niezmienniczego_jest_otwarty_zawarty_w_iteracjach} istnieje niepusty zbiór otwarty $W \subseteq V,$ taki że $W \cup f(W) \cup \ldots \cup f^N(W) \subseteq V$. Na mocy lematu \ref{lemat_3_glownego_artykulu_istnieje_orbita_okresowa_krojaca_sie_z_rodzina_otwartych} istnieje punkt okresowy $x_0$ odwzorowania $f,$ taki że $x_0 \in \textrm{pr}_1(U_i)$ oraz orbita $x_0$ ma niepusty przekrój ze zbiorem $W$. Niech $r > 0$ będzie pierwszą dodatnią liczbą całkowitą dla której $f^r(x_0) \in W$. Wtedy $f^r(x_0), f^{r+1}(x_0), \ldots, f^{r+N-1}(x_0) \in V$. Niech $s \geq 0$ będzie pierwszą nieujemną liczbą całkowitą dla której $f^{r+N+s}(x_0) = x_0$, tj. $r+N+s$ jest okresem punktu $x_0$. Weźmy $y_0,$ takie że $(x_0, y_0) \in U_i$. Ponieważ wszystkie odwzorowania włóknowe $(g_x)$ odwzorowania $F$ są ,,na'', to istnieje punkt $y^* \in (0,1),$ taki że $F^s(f^{r+N}(x_0), y^*) = (x_0, y_0)$. Uzasadnienie:


\begin{equation}
\begin{aligned}
F^s(f^{r+N}(x_0), y^*) & = \Big( f^{r+N+s}(x_0), \, g_{x_s}\big(g_{x_{s-1}}\big( \ldots \big(g_{x_1}(y^*)\big) \ldots \big)\big) \Big) \\
& = \Big( x_0, \, \big(g_{x_s} \circ g_{x_{s-1}} \circ \ldots \circ g_{x_1} \big)(y^*) \Big),
\end{aligned}
\end{equation}
gdzie $x_1 = f^{r+N}(x_0), \, x_{k+1} = f(x_k)$ dla $k=1,2,\ldots,s-1.$
Oznaczmy  $\big(g_{x_s} \circ g_{x_{s-1}} \circ \ldots \circ g_{x_1} \big)$ przez $g_s.$ Funkcja $g_s$ jest ,,na'' jako złożenie funkcji ,,na''. W związku z tym, dla każdego $z \in [0,1]$ istnieje $z^* \in [0,1],$ takie że $g_s(z^*) = z.$ Ponadto jeżeli $z \in (0,1)$, to $z^* \in (0,1)$ bo $g_s$ pozostawia krańce przedziału $[0,1]$ niezmienione.



Przypadek 1. $z = \textrm{pr}_2(F^r(x_0, y_0))$ jest różne od 0 i 1.
Oznaczmy przez $g$ odwzorowanie z $C(I)$ posiadające następujące trzy własności:
\begin{enumerate}[label=(g\arabic*)]
\item \label{g_jeden} $d_1(g, \mathrm{Id}) < \frac{\epsilon}{4}$,
\item \label{g_dwa} $g$ jest odwzorowaniem niemalejącym, pozostawiającym końce przedziału $I$ niezmienione,
\item \label{g_trzy} $g^N(z) = y^*$.
\end{enumerate}

Następnie, rozważmy odwzorowanie $h \in C(I)$ posiadające trzy następujące własności:
\begin{enumerate}[label=(h\arabic*)]
\item \label{h_jeden} $d_1(h, g_{x_0}) < \frac{\epsilon}{4}$,
\item \label{h_dwa} $h(y_0) = g_{x_0}(y_0)$,
\item \label{h_trzy} $h$ jest stałe na zwartym odcinku $[a,b] \subseteq \textrm{pr}_2(U_i)$ zawierającym punkt $y_0$ w swoim wnętrzu.
\end{enumerate}

Weźmy teraz odwzorowanie $G = (f, \widetilde{g}_x) \in \mathcal{F},$ takie że $d_2(G, F) < \frac{\epsilon}{2}$ oraz

\begin{equation}
\label{rownanie_g_z_falka}
    \widetilde{g}_x =
    \begin{cases}
        h & \text{if $x=x_0$,}\\
        g & \text{if $x \in \{f^k(x_0) : r \leq k \leq r+N-1\}$,}\\
        g_x & \text{if $x \in \{f^k(x_0) : 1 \leq k \leq r-1$ or $r+N \leq k \leq r+N+s-1\}$.}
    \end{cases}
\end{equation}

Odwzorowanie takie istnieje na mocy lematu \ref{lemat_4_glownego_artykulu}, ponieważ
$$d_1(\widetilde{g}_{x_0}, g_{x_0}) = d_1(h, g_{x_0}) < \frac{\epsilon}{4}$$
oraz dla $x \in \{f^r(x_0), f^{r+1}(x_0), \ldots, f^{r+N-1}(x_0)\}$,
$$d_1(\widetilde{g}_x, g_x) = d_1(g, g_x) \leq d_1(g, \mathrm{Id}) + d_1(\mathrm{Id}, g_x) < \frac{\epsilon}{4} + \frac{\epsilon}{4} = \frac{\epsilon}{2}$$.

Punkt $(x_0, y_0) \in U_i$ jest punktem okresowym odwzorowania $G$, ponieważ

\begin{equation}
\begin{split}
G^{r+N+s}(x_0, y_0) & \stackrel{(1)}{=} G^{r+N+s-1}(G(x_0, y_0)) \\
& \stackrel{(2)}{=} G^{r+N+s-1}(F(x_0, y_0)) \stackrel{(3)}{=} G^{N+s}(F^r(x_0, y_0)) \\
& \stackrel{(4)}{=} G^{N+s}(f^r(x_0), z) \stackrel{(5)}{=} G^s(f^{r+N}(x_0), y^*) \\
& \stackrel{(6)}{=} F^s(f^{r+N}(x_0), y^*) \stackrel{(7)}{=} (x_0, y_0).
\end{split}
\end{equation}

Szczegółowe uzasadnienie powyższych równości:
\begin{enumerate}
\item Nie musimy rozważać odwracalności $G$ gdyż $r > 0, \, N > 0, \, s \geq 0$, a więc $r+N+s \geq 2.$
\item $G = (f, \widetilde{g}_x), \, \widetilde{g}_{x_0} = h$ a $h(y_0) = g_{x_0}(y_0).$
\item Z równania \ref{rownanie_g_z_falka} mamy $\widetilde{g}_x = g_x,$ dla $x \in \{f^k(x_0) : 1 \leq k \leq r-1\}.$
\item Wynika z definicji $z.$
\item Z równania \ref{rownanie_g_z_falka} dla $x \in \{f^k(x_0) : r \leq k \leq r+N-1\},$ oraz \ref{g_trzy}.
\item Z równania \ref{rownanie_g_z_falka} mamy $\widetilde{g}_x = g_x,$ dla $x \in \{f^k(x_0) : r+N \leq k \leq r+N+s-1\}.$
\item Wprost z definicji $y^*.$
\end{enumerate}

Ponadto, ze względu na \ref{h_trzy} zachodzi

\begin{equation}
\begin{aligned}
G^{r+N+s}(\{x_0\} \times [a, b]) & = G^{r+N+s-1}\big(G(\{x_0\} \times [a, b])\big) \\
& = G^{r+N+s-1}\big(\{f(x_0)\} \times \widetilde{g}_x([a,b])\big) \\
& \stackrel{\ref{h_trzy}}{=} G^{r+N+s-1}\big(\{f(x_0)\} \times \widetilde{g}_{x_0}([a,b])\big) \\
& \stackrel{\ref{h_dwa}}{=} G^{r+N+s-1}\big(\{f(x_0)\} \times \{g_{x_0}(y_0)\}\big) \\
& = G^{r+N+s-1}\big(\{(f(x_0), g_{x_0}(y_0))\}\big) \\
& = G^{r+N+s-1}\big(\{F(x_0, y_0)\}\big) \\
& = \{(x_0, y_0)\}.
\end{aligned}
\end{equation}


Z faktu $y_0 \in (a, b)$ wynika, że każde odwzorowanie $\widetilde{G} \in \mathcal{F}$ dostatecznie bliskie $G$ posiada własność
$$\widetilde{G}^{r+N+s}(\{x_0\} \times [a, b]) \subseteq \{x_0\} \times [a, b].$$

$\widetilde G^{r+N+s}$ na pierwszej współrzędnej przeprowadza $x_0$ na $x_0$, zaś na drugiej przeprowadza odcinek $[a,b]$ w $[a,b]$. Każde ciągłe odwzorowanie odcinka w siebie ma punkt stały, a punkt stały iteraty jest punktem okresowym oryginalnego odwzorowania.
Zatem $\widetilde{G}$ posiada punkt okresowy w $\{x_0\} \times [a, b] \subseteq U_i$.

Zatem $G \in \mathcal{F}_{SO}^i$, co kończy dowód dla przypadku 1.

Przypadek 2. $\textrm{pr}_2(F^r(x_0, y_0))$ jest równy $0$ lub $1$.

W takim przypadku użyjemy lematu \ref{lemat_4_glownego_artykulu} aby dostać odwzorowanie $H = (f, h_x) \in \mathcal{F},$ takie że $d_2(H, F) < \frac{\epsilon}{2}$ i dla $x \in \{x_0, f(x_0), \ldots, f^{r-1}(x_0)\}$ odwzorowania włóknowe $h_x$ są ściśle rosnące. Takie odwzorowanie istnieje, ponieważ każdą funkcję niemalejącą możemy jednostajnie przybliżać funkcjami ściśle rosnącymi.

 
Ponieważ $y_0$ jest różnie od $0$ i $1$ dostajemy, że $\textrm{pr}_2(H^r(x_0, y_0))$ również jest różne od $0$~i~$1$. (Ponieważ $H = (f, h_x) \in \mathcal{F}$, to każde z odwzorowań $h_x$ pozostawia niezmienione punkty $0$ i $1$. Wprowadźmy oznaczenie $\textrm{pr}_2(H^r(x_0, y_0)) = (h_{f^{r-1}(x_0)} \circ \ldots \circ h_{f(x_0)} \circ h_{x_0})(y_0) = h_0(y_0)$. Odwzorowanie $h_0$ jako złożenie odwzorowań ściśle rosnących i pozostawiających punkty $0$ i $1$ niezmienione jest ściśle rosnące i pozostawia niezmienione punkty $0$ i $1$.  Ponieważ $y_0 > 0$ to oczywiście $h_0(y_0) > h_0(0) = 0$. Natomiast ze względu na fakt, że $y_0 < 1$, to zachodzi $h_0(y_0) < h_0(1) = 1$.)
 
Następnie korzystając z przypadku 1 dostajemy odwzorowanie $G \in \mathcal{F}_{SO}^i$, dla którego zachodzi nierówność $d_2(G, H) < \frac{\epsilon}{2}$. Wówczas, korzystając z nierówności trójkąta otrzymujemy $d_2(G, F) \leq d_2(G, H) + d_2(H, F) < \frac{\epsilon}{2} + \frac{\epsilon}{2} = \epsilon$.


\end{proof}




% END MOJE ====================================================



{\backmatter \chapter{Podsumowanie}}
Podsumowanie w pracach matematycznych nie jest obligatoryjne. Warto jednak na zakończenie krótko napisać, co udało nam się zrobić w pracy, a czasem także o tym, czego nie udało się zrobić.

{\backmatter \chapter{Dodatek}}
Dodatek w pracach matematycznych również nie jest wymagany. Można w nim przedstawić np. jakiś dłuższy dowód, który z pewnych przyczyn pominęliśmy we właściwej części pracy lub (np. w przypadku prac statystycznych) umieścić dane, które analizowaliśmy.

%%%%%%%%%%%%%%%%%%%%%%%%%%%%%%%%%%%%%%%%%%%%%%%%%%%%%%%%%
% BIBLIOGRAFIA
% W tworzeniu bibliografii najlepiej korzystać z BibTex'a, 
% który jest częścią systemu Tex. W naszym przypadku funkcję 
% przechowalni literatury, do której się odwołujemy, pełni 
% plik bibliografia.bib. Nie musimy ręcznie dodawać nowych 
% pozycji do bibliografii. Możemy wejść np. na stronę 
% https://mathscinet.ams.org/mathscinet/index.html, 
% znaleźć odpowiednią pozycję, wybrać ją, a następnie zmienić 
% 'Select alternative format' na BibTeX, skopiować uzyskany 
% tekst, wkleić do pliku bibliografia.bib i skompilować. 
% Gotowe informacje do pliku bibliografia.bib można znaleźć 
% także na https://arxiv.org - gdy znajdziemy interesującą nas 
% pracę, szukamy 'References & Citations' i klikamy 'NASA ADS', 
% a potem 'Bibtex entry for this abstract' 
% i postępujemy tak jak wcześniej.
%%%%%%%%%%%%%%%%%%%%%%%%%%%%%%%%%%%%%%%%%%%%%%%%%%%%%%%%%
\newpage
% w nawiasie klamrowym wpisujemy nazwę pliku z bibliografią w formacie .bib
\bibliografia{bibliografia.bib} 
\end{document}